\documentclass[UTF8, 12pt, fontset=fandol, oneside]{ctexart}
\usepackage{researchnotes}

\title{第 2 章\quad 凸集}
\author{}
\date{}

\begin{document}

\maketitle

\section*{2.1\quad 仿射集合和凸集}

\subsection*{2.1.1\quad 直线与线段}

设 $x_1\ne x_2$ 为 $\mathbb{R}^n$ 空间中的两个点，那么具有下列形式的点
\[
  y=\theta x_1+(1-\theta)x_2,\quad \theta\in\mathbb{R}
\]
组成一条穿越 $x_1$ 和 $x_2$ 的直线。参数 $\theta=0$ 对应 $y=x_2$，而 $\theta=1$ 表明 $y=x_1$。参数 $\theta$ 的值在 0 和 1 之间变动，构成了 $x_1$ 和 $x_2$ 之间的（闭）线段。

$y$ 的表示形式
\[
  y=x_2+\theta(x_1-x_2)
\]
给出了另一种解释：$y$ 是基点 $x_2$（对应 $\theta=0$）和方向 $x_1-x_2$（由 $x_2$ 指向 $x_1$）乘以参数 $\theta$ 的和。因此，$\theta$ 给出了 $y$ 在由 $x_2$ 通向 $x_1$ 的路上的位置。当 $\theta$ 由 0 增加到 1，点 $y$ 相应地由 $x_2$ 移动到 $x_1$。如果 $\theta>1$，点 $y$ 在超越了 $x_1$ 的直线上。图2.1给出了直观的解释。

\subsection*{2.1.2\quad 仿射集合}

如果通过集合 $C\subseteq\mathbb{R}^n$ 中任意两个不同点的直线仍然在集合 $C$ 中，那么称集合 $C$ 是\textbf{仿射的}。也就是说，$C\subseteq\mathbb{R}^n$ 是仿射的等价于：对于任意 $x_1,x_2\in C$ 及 $\theta\in\mathbb{R}$ 有
\[
  \theta x_1+(1-\theta)x_2\in C.
\]
换而言之，$C$ 包含了 $C$ 中任意两点的系数之和为 1 的线性组合。

这个概念可以扩展到多个点的情况。如果 $\theta_1+\cdots+\theta_k=1$，我们称具有 $\theta_1x_1+\cdots+\theta_kx_k$ 形式的点为 $x_1,\cdots,x_k$ 的\textbf{仿射组合}。利用仿射集合的定义（即仿射集合包含其中任意两点的仿射组合），我们可以归纳出以下结论：一个仿射集合包含其中任意点的仿射组合，即如果 $C$ 是一个仿射集合，$x_1,\cdots,x_k\in C$，并且 $\theta_1+\cdots+\theta_k=1$，那么 $\theta_1x_1+\cdots+\theta_kx_k$ 仍然在 $C$ 中。

如果 $C$ 是一个仿射集合并且 $x_0\in C$，则集合
\[
  V=C-x_0=\{x-x_0\mid x\in C\}
\]
是一个子空间，即关于加法和数乘是封闭的。为说明这一点，设 $v_1,v_2\in V$，$\alpha,\beta\in\mathbb{R}$，则有 $v_1+x_0\in C$，$v_2+x_0\in C$。因为 $C$ 是仿射的，且 $\alpha+\beta+(1-\alpha-\beta)=1$，所以
\[
  \alpha v_1+\beta v_2+x_0
  =\alpha(v_1+x_0)+\beta(v_2+x_0)+(1-\alpha-\beta)x_0\in C,
\]
由 $\alpha v_1+\beta v_2+x_0\in C$，我们可知 $\alpha v_1+\beta v_2\in V$。

因此，仿射集合 $C$ 可以表示为
\[
  C=V+x_0=\{v+x_0\mid v\in V\},
\]
即一个子空间加上一个偏移。与仿射集合 $C$ 相关联的子空间 $V$ 与 $x_0$ 的选取无关，所以 $x_0$ 可以是 $C$ 中的任意一点。我们定义仿射集合 $C$ 的维数为子空间 $V=C-x_0$ 的维数，其中 $x_0$ 是 $C$ 中的任意元素。

\noindent\textbf{例 2.1}\quad 线性方程组的解集。线性方程组的解集 $C=\{x\mid Ax=b\}$，其中 $A\in\mathbb{R}^{m\times n}$，$b\in\mathbb{R}^m$，是一个仿射集合。为说明这点，设 $x_1,x_2\in C$，即 $Ax_1=b$，$Ax_2=b$。则对于任意 $\theta$，我们有
\[
\begin{aligned}
  A\bigl(\theta x_1+(1-\theta)x_2\bigr)
  &=\theta Ax_1+(1-\theta)Ax_2\\
  &=\theta b+(1-\theta)b\\
  &=b,
\end{aligned}
\]
这表明任意的仿射组合 $\theta x_1+(1-\theta)x_2$ 也在 $C$ 中，并且与仿射集合 $C$ 相关联的子空间就是 $A$ 的零空间。

反之任意仿射集合可以表示为一个线性方程组的解集。

我们称由集合 $C\subseteq\mathbb{R}^n$ 中的点的所有仿射组合组成的集合为 $C$ 的\textbf{仿射包}，记为 $\mathbf{aff}\,C$：
\[
  \mathbf{aff}\,C=\{\theta_1x_1+\cdots+\theta_kx_k\mid x_1,\cdots,x_k\in C,\ \theta_1+\cdots+\theta_k=1\}.
\]
仿射包是包含 $C$ 的最小的仿射集合，也就是说：如果 $S$ 是满足 $C\subseteq S$ 的仿射集合，那么 $\mathbf{aff}\,C\subseteq S$。

\subsection*{2.1.3\quad 仿射维数与相对内部}

我们定义集合 $C$ 的\textbf{仿射维数}为其仿射包的维数。仿射维数在凸分析及凸优化中十分有用，但它与其他维数的定义常常不相容。作为一个例子，考虑 $\mathbb{R}^2$ 上的单位圆环
\[
  \{x\in\mathbb{R}^2\mid x_1^2+x_2^2=1\}.
\]
它的仿射包是全空间 $\mathbb{R}^2$，所以其仿射维数为 2。但是，在其他大多数维数的定义下，$\mathbb{R}^2$ 上的单位圆环的维数为 1。

如果集合 $C\subseteq\mathbb{R}^n$ 的仿射维数小于 $n$，那么这个集合在仿射集合 $\mathbf{aff}\,C\ne\mathbb{R}^n$ 中。我们定义集合 $C$ 的\textbf{相对内部}为 $\mathbf{aff}\,C$ 的内部，记为 $\mathbf{relint}\,C$，即
\[
  \mathbf{relint}\,C=\{x\in C\mid B(x,r)\cap \mathbf{aff}\,C\subseteq C\text{ 对于某些 }r>0\},
\]
其中 $B(x,r)=\{y\mid \|y-x\|\le r\}$，即半径为 $r$、中心为 $x$ 并由范数 $\|\cdot\|$ 定义的球（这里的 $\|\cdot\|$ 可以是任意范数，并且所有范数定义了相同的相对内部）。我们于是可以定义集合 $C$ 的\textbf{相对边界}为 $\mathbf{cl}\,C\setminus\mathbf{relint}\,C$，此处 $\mathbf{cl}\,C$ 表示 $C$ 的闭包。

\noindent\textbf{例 2.2}\quad 考虑 $\mathbb{R}^3$ 中处于 $(x_1,x_2)$ 平面的一个正方形，定义
\[
  C=\{x\in\mathbb{R}^3\mid -1\le x_1\le 1,\ -1\le x_2\le 1,\ x_3=0\}.
\]
其仿射包为 $(x_1,x_2)$-平面，即 $\mathbf{aff}\,C=\{x\in\mathbb{R}^3\mid x_3=0\}$。$C$ 的内部为空，但其相对内部为
\[
  \mathbf{relint}\,C=\{x\in\mathbb{R}^3\mid -1<x_1<1,\ -1<x_2<1,\ x_3=0\}.
\]
$C$（在 $\mathbb{R}^3$ 中）的边界是其自身，而相对边界是其边框：
\[
  \{x\in\mathbb{R}^3\mid \max\{|x_1|,|x_2|\}=1,\ x_3=0\}.
\]

\subsection*{2.1.4\quad 凸集}

集合 $C$ 被称为\textbf{凸集}，如果 $C$ 中任意两点间的线段仍然在 $C$ 中，即对于任意 $x_1,x_2\in C$ 和满足 $0\le\theta\le 1$ 的 $\theta$ 都有
\[
  \theta x_1+(1-\theta)x_2\in C.
\]
粗略地，如果集合中的每一点都可以被其他点沿着它们之间一条无阻碍的路径看见，那么这个集合就是凸集。所谓无阻碍，是指整条路径都在集合中。由于仿射集包含穿过集合中任意不同两点的整条直线，任意不同两点间的线段自然也在集合中。因而仿射集是凸集。图2.2显示了 $\mathbb{R}^2$ 空间中一些简单的凸和非凸集合。

我们称点 $\theta_1x_1+\cdots+\theta_kx_k$ 为点 $x_1,\cdots,x_k$ 的一个\textbf{凸组合}，其中 $\theta_1+\cdots+\theta_k=1$ 并且 $\theta_i\ge 0$，$i=1,\cdots,k$。与仿射集合类似，一个集合是凸集等价于集合包含其中所有点的凸组合。点的凸组合可以看做它们的混合或加权平均，$\theta_i$ 代表混合时 $x_i$ 所占的份数。

我们称集合 $C$ 中所有点的凸组合的集合为其\textbf{凸包}，记为 $\mathbf{conv}\,C$：
\[
  \mathbf{conv}\,C=\{\theta_1x_1+\cdots+\theta_kx_k\mid x_i\in C,\ \theta_i\ge 0,\ i=1,\cdots,k,\ \theta_1+\cdots+\theta_k=1\}.
\]
顾名思义，凸包 $\mathbf{conv}\,C$ 总是凸的。它是包含 $C$ 的最小的凸集。也就是说，如果 $B$ 是包含 $C$ 的凸集，那么 $\mathbf{conv}\,C\subseteq B$。图2.3显示了凸包的定义。

凸组合的概念可以扩展到无穷级数、积分以及大多数形式的概率分布。假设 $\theta_1,\theta_2,\cdots$ 满足
\[
  \theta_i\ge 0,\quad i=1,2,\cdots,\qquad
  \sum_{i=1}^{\infty}\theta_i=1,
\]
并且 $x_1,x_2,\cdots\in C$，其中 $C\subseteq\mathbb{R}^n$ 为凸集。那么，如果下面的级数收敛，我们有
\[
  \sum_{i=1}^{\infty}\theta_i x_i\in C.
\]
更一般地，假设 $p:\mathbb{R}^n\to\mathbb{R}$ 对所有 $x\in C$ 满足 $p(x)\ge 0$，并且 $\int_C p(x)\,dx=1$，其中 $C\subseteq\mathbb{R}^n$ 是凸集。那么，如果下面的积分存在，我们有
\[
  \int_C p(x)x\,dx\in C,
\]
最一般的情况，设 $C\subseteq\mathbb{R}^n$ 是凸集，$x$ 是随机变量，并且 $x\in C$ 的概率为 1，那么 $\mathbf{E}x\in C$。事实上，这一形式包含了前述的特殊情况。例如，假设随机变量 $x$ 只在 $x_1$ 和 $x_2$ 中取值，其概率分别为 $\operatorname{prob}(x=x_1)=\theta$ 和 $\operatorname{prob}(x=x_2)=1-\theta$，其中 $0\le\theta\le 1$。于是，$\mathbf{E}x=\theta x_1+(1-\theta)x_2$，即回到了两个点的简单的凸组合。

\subsection*{2.1.5\quad 锥}

如果对于任意 $x\in C$ 和 $\theta\ge 0$ 都有 $\theta x\in C$，我们称集合 $C$ 是\textbf{锥}或者\textbf{非负齐次}。如果集合 $C$ 是锥，并且是凸的，则称 $C$ 为\textbf{凸锥}，即对于任意 $x_1,x_2\in C$ 和 $\theta_1,\theta_2\ge 0$，都有
\[
  \theta_1x_1+\theta_2x_2\in C.
\]
在几何上，具有此类形式的点构成了二维的扇形，这个扇形以 0 为顶点，边通过 $x_1$ 和 $x_2$（如图2.4所示）。

具有 $\theta_1x_1+\cdots+\theta_kx_k$，$\theta_1,\cdots,\theta_k\ge 0$ 形式的点称为 $x_1,\cdots,x_k$ 的\textbf{锥组合}（或非负线性组合）。如果 $x_i$ 均属于凸锥 $C$，那么，$x_i$ 的每一个锥组合也在 $C$ 中。反言之，集合 $C$ 是凸锥的充要条件是它包含其元素的所有锥组合。如同凸（或仿射）组合一样，锥组合的概念可以扩展到无穷级数和积分中。

集合 $C$ 的\textbf{锥包}是 $C$ 中元素的所有锥组合的集合，即
\[
  \{\theta_1x_1+\cdots+\theta_kx_k\mid x_i\in C,\ \theta_i\ge 0,\ i=1,\cdots,k\},
\]
它是包含 $C$ 的最小的凸锥（如图2.5所示）。

\section*{2.2\quad 重要的例子}

本节将描述一些重要的凸集，这些凸集在本书的后续部分将多次遇见。首先介绍一些简单的例子。

\begin{itemize}
  \item 空集 $\emptyset$、任意一个点（即单点集）$\{x_0\}$、全空间 $\mathbb{R}^n$ 都是 $\mathbb{R}^n$ 的仿射（自然也是凸的）子集。
  \item 任意直线是仿射的。如果直线通过零点，则是子空间，因此，也是凸锥。
  \item 一条线段是凸的，但不是仿射的（除非退化为一个点）。
  \item 一条射线，即具有形式 $\{x_0+\theta v\mid \theta\ge 0\}$，$v\ne 0$ 的集合，是凸的，但不是仿射的。如果射线的基点 $x_0$ 是 0，则它是凸锥。
  \item 任意子空间是仿射的、凸锥（自然是凸的）。
\end{itemize}

\subsection*{2.2.1\quad 超平面与半空间}

超平面是具有下面形式的集合
\[
  \{x\mid a^Tx=b\},
\]
其中 $a\in\mathbb{R}^n$，$a\ne 0$ 且 $b\in\mathbb{R}$。解析地，超平面是关于 $x$ 的非平凡线性方程的解空间（因此是一个仿射集合）。几何上，超平面 $\{x\mid a^Tx=b\}$ 可以解释为与给定向量 $a$ 的内积为常数的点的集合；也可以看成法线方向为 $a$ 的超平面，而常数 $b\in\mathbb{R}$ 决定了这个平面从原点的偏移。为更好地理解这个几何解释，可以将超平面表示成下面的形式
\[
  \{x\mid a^T(x-x_0)=0\},
\]
其中 $x_0$ 是超平面上的任意一点（即任意满足 $a^Tx_0=b$ 的点）。进一步，可以表示为
\[
  \{x\mid a^T(x-x_0)=0\}=x_0+a^\perp,
\]
这里 $a^\perp$ 表示 $a$ 的正交补，即与 $a$ 正交的向量的集合：
\[
  a^\perp=\{v\mid a^Tv=0\}.
\]
从中可以看出，超平面由偏移 $x_0$ 加上所有正交于（法）向量 $a$ 的向量构成。这些几何解释可见图2.6。

一个超平面将 $\mathbb{R}^n$ 划分为两个半空间。（闭的）半空间是具有下列形式的集合：
\begin{equation}
  \{x\mid a^Tx\le b\},
  \tag{2.1}
\end{equation}
即（非平凡的）线性不等式的解空间，其中 $a\ne 0$。半空间是凸的，但不是仿射的，如图2.7所示。

半空间 (2.1) 也可表示为
\begin{equation}
  \{x\mid a^T(x-x_0)\le 0\},
  \tag{2.2}
\end{equation}
其中 $x_0$ 是相应超平面上的任意一点，即 $x_0$ 满足 $a^Tx_0=b$。表达式 (2.2) 有一个简单的几何解释：半空间由 $x_0$ 加上任意与（向外的法）向量 $a$ 呈钝角（或直角）的向量组成，如图2.8所示。

半空间 (2.1) 的边界是超平面 $\{x\mid a^Tx=b\}$。集合 $\{x\mid a^Tx<b\}$ 是半空间 $\{x\mid a^Tx\le b\}$ 的内部，称为\textbf{开半空间}。

\subsection*{2.2.2\quad Euclid 球和椭球}

$\mathbb{R}^n$ 中的 Euclid 球（或简称为球）具有下面的形式，
\[
  B(x_c,r)=\{x\mid \|x-x_c\|_2\le r\}
  =\{x\mid (x-x_c)^T(x-x_c)\le r^2\},
\]
其中 $r>0$，$\|\cdot\|_2$ 表示 Euclid 范数，即 $\|u\|_2=(u^Tu)^{1/2}$。向量 $x_c$ 是球心，标量 $r$ 为半径。$B(x_c,r)$ 由距离球心 $x_c$ 距离不超过 $r$ 的所有点组成。Euclid 球的另一个常见的表达式为，
\[
  B(x_c,r)=\{x_c+ru\mid \|u\|_2\le 1\}.
\]
Euclid 球是凸集，即如果 $\|x_1-x_c\|_2\le r$，$\|x_2-x_c\|_2\le r$，并且 $0\le\theta\le 1$，那么
\[
\begin{aligned}
  \|\theta x_1+(1-\theta)x_2-x_c\|_2
  &=\|\theta(x_1-x_c)+(1-\theta)(x_2-x_c)\|_2\\
  &\le \theta\|x_1-x_c\|_2+(1-\theta)\|x_2-x_c\|_2\\
  &\le r.
\end{aligned}
\]
（此处利用了 $\|\cdot\|_2$ 的齐次性和三角不等式，参见 §A.1.2。）

一类相关的凸集是\textbf{椭球}，它们具有如下的形式，
\begin{equation}
  \mathcal{E}=\{x\mid (x-x_c)^TP^{-1}(x-x_c)\le 1\},
  \tag{2.3}
\end{equation}
其中 $P=P^T\succ 0$，即 $P$ 是对称正定矩阵。向量 $x_c\in\mathbb{R}^n$ 为椭球的中心。矩阵 $P$ 决定了椭球从 $x_c$ 向各个方向扩展的幅度。$\mathcal{E}$ 的半轴长度由 $\sqrt{\lambda_i}$ 给出，这里 $\lambda_i$ 为 $P$ 的特征值。球可以看成 $P=r^2I$ 的椭球。图2.9显示了 $\mathbb{R}^2$ 上的一个椭球。

椭球另一个常用的表示形式是
\begin{equation}
  \mathcal{E}=\{x_c+Au\mid \|u\|_2\le 1\},
  \tag{2.4}
\end{equation}
其中 $A$ 是非奇异的方阵。在此类表示形式中，我们可以不失一般性地假设 $A$ 对称正定。取 $A=P^{1/2}$，这个表达式给出了由式 (2.3) 定义的椭球。当式 (2.4) 中的矩阵 $A$ 为对称半正定矩阵，但奇异时，集合 (2.4) 称为\textbf{退化的椭球}，其仿射维数等于 $A$ 的秩。退化的椭球也是凸的。

\subsection*{2.2.3\quad 范数球和范数锥}

设 $\|\cdot\|$ 是 $\mathbb{R}^n$ 中的范数（参见附录 A.1.2）。由范数的一般性质可知，以 $r$ 为半径，$x_c$ 为球心的范数球 $\{x\mid \|x-x_c\|\le r\}$ 是凸的。关于范数 $\|\cdot\|$ 的范数锥是集合
\[
  C=\{(x,t)\mid \|x\|\le t\}\subseteq\mathbb{R}^{n+1}.
\]
顾名思义，它是一个凸锥。

\noindent\textbf{例 2.3}\quad 二阶锥是由 Euclid 范数定义的范数锥，即
\[
  C=\{(x,t)\in\mathbb{R}^{n+1}\mid \|x\|_2\le t\}
  =
  \left\{
  \begin{bmatrix} x\\ t \end{bmatrix}
  \ \middle|\
  \begin{bmatrix} x\\ t \end{bmatrix}^{T}
  \begin{bmatrix} I & 0\\ 0 & -1 \end{bmatrix}
  \begin{bmatrix} x\\ t \end{bmatrix}
  \le 0,\ t\ge 0
  \right\}.
\]
二阶锥的其他名字也常常被使用。它由二次不等式定义，因此也被称为二次锥。同时，也称其为 Lorentz 锥或冰激凌锥。图2.10显示了 $\mathbb{R}^3$ 上一个的二阶锥。

\subsection*{2.2.4\quad 多面体}

多面体被定义为有限个线性等式和不等式的解集，
\begin{equation}
  \mathcal{P}=\{x\mid a_j^Tx\le b_j,\ j=1,\cdots,m,\ c_j^Tx=d_j,\ j=1,\cdots,p\}.
  \tag{2.5}
\end{equation}
因此，多面体是有限个半空间和超平面的交集。仿射集合（例如子空间、超平面、直线）、射线、线段和半空间都是多面体。显而易见，多面体是凸集。有界的多面体有时也称为\textbf{多胞形}，但也有一些作者反过来使用这两个概念（即用多胞形表示具有 (2.5) 形式的集合，而当其有界时称为多面体）。图 2.11 显示了一个由五个半空间的交集定义的多面体。

可以方便地使用紧凑表达式
\begin{equation}
  \mathcal{P}=\{x\mid Ax\preceq b,\ Cx=d\}
  \tag{2.6}
\end{equation}
来表示 (2.5)，其中，
\[
  A=
  \begin{bmatrix}
    a_1^T\\
    \vdots\\
    a_m^T
  \end{bmatrix},
  \qquad
  C=
  \begin{bmatrix}
    c_1^T\\
    \vdots\\
    c_p^T
  \end{bmatrix},
\]
此处的 $\preceq$ 代表 $\mathbb{R}^m$ 上的向量不等式或分量不等式；$u\preceq v$ 表示 $u_i\le v_i$，$i=1,\cdots,m$。

\noindent\textbf{例 2.4}\quad 非负象限是具有非负分量的点的集合，即
\[
  \mathbb{R}_+^n=\{x\in\mathbb{R}^n\mid x_i\ge 0,\ i=1,\cdots,n\}
  =\{x\in\mathbb{R}^n\mid x\succeq 0\}.
\]
（此处 $\mathbb{R}_+$ 表示非负实数的集合，即 $\mathbb{R}_+=\{x\in\mathbb{R}\mid x\ge 0\}$。）非负象限既是多面体也是锥（因此称为\textbf{多面体锥}）。

\noindent\textbf{单纯形}

单纯形是一类重要的多面体。设 $k+1$ 个点 $v_0,\cdots,v_k\in\mathbb{R}^n$ 仿射独立，即 $v_1-v_0,\cdots,v_k-v_0$ 线性独立，那么，这些点决定了一个单纯形，如下
\begin{equation}
  C=\mathbf{conv}\{v_0,\cdots,v_k\}
  =\{\theta_0v_0+\cdots+\theta_kv_k\mid \theta\succeq 0,\ \mathbf{1}^T\theta=1\},
  \tag{2.7}
\end{equation}
其中 $\mathbf{1}$ 表示所有分量均为一的向量。这个单纯形的仿射维数为 $k$，因此也称为 $\mathbb{R}^n$ 空间的 $k$ 维单纯形。

\noindent\textbf{例 2.5}\quad 一些常见的单纯形。1 维单纯形是一条线段；2 维单纯形是一个三角形（包含其内部）；3 维单纯形是一个四面体。

单位单纯形是由零向量和单位向量 $0,e_1,\cdots,e_n\in\mathbb{R}^n$ 决定的 $n$ 维单纯形。它可以表示为满足下列条件的向量的集合，
\[
  x\succeq 0,\qquad \mathbf{1}^Tx\le 1.
\]
概率单纯形是由单位向量 $e_1,\cdots,e_n\in\mathbb{R}^n$ 决定的 $n-1$ 维单纯形。它是满足下列条件的向量的集合，
\[
  x\succeq 0,\qquad \mathbf{1}^Tx=1.
\]
概率单纯形中的向量对应于含有 $n$ 个元素的集合的概率分布，$x_i$ 可理解为第 $i$ 个元素的概率。

为用多面体来描述单纯形 (2.7)，我们采用以下步骤将其变换为 (2.6) 的形式。由定义可知，$x\in C$ 的充要条件是，对于某些 $\theta\succeq 0,\ \mathbf{1}^T\theta=1$，有 $x=\theta_0v_0+\theta_1v_1+\cdots+\theta_kv_k$。等价地，如果定义 $y=(\theta_1,\cdots,\theta_k)$ 和
\[
  B=\begin{bmatrix} v_1-v_0 & \cdots & v_k-v_0 \end{bmatrix}\in\mathbb{R}^{n\times k},
\]
我们知道 $x\in C$ 的充要条件是
\begin{equation}
  x=v_0+By
  \tag{2.8}
\end{equation}
对于 $y\succeq 0,\ \mathbf{1}^Ty\le 1$ 成立。注意到 $v_0,\cdots,v_k$ 仿射独立意味着矩阵 $B$ 的秩为 $k$。因此，存在非奇异矩阵 $A=(A_1,A_2)\in\mathbb{R}^{n\times n}$ 使得
\[
  AB=
  \begin{bmatrix}
    A_1\\
    A_2
  \end{bmatrix}B
  =
  \begin{bmatrix}
    I\\
    0
  \end{bmatrix}.
\]
用 $A$ 左乘 (2.8)，我们得到
\[
  A_1x=A_1v_0+y,\qquad A_2x=A_2v_0.
\]
从中可以看出，$x\in C$ 当且仅当 $A_2x=A_2v_0$ 并且向量 $y=A_1x-A_1v_0$ 满足 $y\succeq 0$ 和 $\mathbf{1}^Ty\le 1$。换言之我们得到了 $x\in C$ 的充要条件，
\[
  A_2x=A_2v_0,\qquad A_1x\succeq A_1v_0,\qquad \mathbf{1}^TA_1x\le 1+\mathbf{1}^TA_1v_0.
\]
这些是 $x$ 的线性等式和不等式，因此，描述了一个多面体。

\noindent\textbf{多面体的凸包描述}

有限集合 $\{v_1,\cdots,v_k\}$ 的凸包是
\[
  \mathbf{conv}\{v_1,\cdots,v_k\}
  =\{\theta_1v_1+\cdots+\theta_kv_k\mid \theta\succeq 0,\ \mathbf{1}^T\theta=1\}.
\]
它是一个有界的多面体，但是（除非是例如单纯形这样的特殊情况）无法简单地用形如 (2.5) 的式子，即用线性等式和不等式的集合将其表示。

凸包表达式的一个扩展表示是：
\begin{equation}
  \{\theta_1v_1+\cdots+\theta_kv_k\mid \theta_1+\cdots+\theta_m=1,\ \theta_i\ge 0,\ i=1,\cdots,k\},
  \tag{2.9}
\end{equation}
其中 $m\le k$。此处我们考虑 $v_i$ 的非负线性组合，但是仅仅要求前 $m$ 个系数之和为一。此外，我们可以将 (2.9) 解释为点 $v_1,\cdots,v_m$ 的凸包加上点 $v_{m+1},\cdots,v_k$ 的锥包。集合 (2.9) 定义了一个多面体，反之亦然，每个多面体都可以表示为此类形式（虽然这一点此处并未证明）。

``如何表示多面体''这一问题的求解是十分巧妙的，并且有一些非常实用的结果。一个简单的例子是定义在 $\ell_\infty$-范数空间的 $\mathbb{R}^n$ 上的单位球，
\[
  C=\{x\mid |x_i|\le 1,\ i=1,\cdots,n\}.
\]
$C$ 可以由 $2n$ 个线性不等式 $\pm e_i^Tx\le 1$ 表示为 (2.5) 的形式，其中 $e_i$ 表示第 $i$ 维的单位向量。为了将其描述为形如 (2.9) 的凸包，需要至少 $2^n$ 个点：
\[
  C=\mathbf{conv}\{v_1,\cdots,v_{2^n}\},
\]
其中 $v_1,\cdots,v_{2^n}$ 是以 1 和 $-1$ 为分量的全部向量，共 $2^n$ 个。可见，当 $n$ 很大时，这两种描述方式的规模相差极大。

\subsection*{2.2.5\quad 半正定锥}

我们用 $\mathbf{S}^n$ 表示对称 $n\times n$ 矩阵的集合，即
\[
  \mathbf{S}^n=\{X\in\mathbb{R}^{n\times n}\mid X=X^T\},
\]
这是一个维数为 $n(n+1)/2$ 的向量空间。我们用 $\mathbf{S}_+^n$ 表示对称半正定矩阵的集合：
\[
  \mathbf{S}_+^n=\{X\in\mathbf{S}^n\mid X\succeq 0\},
\]
用 $\mathbf{S}_{++}^n$ 表示对称正定矩阵的集合：
\[
  \mathbf{S}_{++}^n=\{X\in\mathbf{S}^n\mid X\succ 0\}.
\]
（这些符号与 $\mathbb{R}_+$ 相对应：$\mathbb{R}_+$ 表示非负实数，而 $\mathbb{R}_{++}$ 表示正实数。）集合 $\mathbf{S}_+^n$ 是一个凸锥：如果 $\theta_1,\theta_2\ge 0$ 并且 $A,B\in\mathbf{S}_+^n$，那么 $\theta_1A+\theta_2B\in\mathbf{S}_+^n$。从半正定矩阵的定义可以直接得到：对于任意 $x\in\mathbb{R}^n$，如果 $A\succeq 0,\ B\succeq 0$ 并且 $\theta_1,\theta_2\ge 0$，那么，我们有
\[
  x^T(\theta_1A+\theta_2B)x=\theta_1x^TAx+\theta_2x^TBx\ge 0.
\]

\noindent\textbf{例 2.6}\quad $\mathbf{S}^2$ 上的半正定锥。我们有
\[
  X=
  \begin{bmatrix}
    x & y\\
    y & z
  \end{bmatrix}\in\mathbf{S}_+^2
  \Longleftrightarrow
  x\ge 0,\quad z\ge 0,\quad xz\ge y^2.
\]
图2.12显示了这个锥的边界，按 $(x,y,z)$ 表示在 $\mathbb{R}^3$ 中。

\section*{2.3\quad 保凸运算}

本节将描述一些保凸运算，利用它们，我们可以从凸集构造出其他凸集。这些运算与 §2.2 中描述的凸集的简单例子一起构成了凸集的演算，可以用来确定或构建集合的凸性。

\subsection*{2.3.1\quad 交集}

交集运算是保凸的：如果 $S_1$ 和 $S_2$ 是凸集，那么 $S_1\cap S_2$ 也是凸集。这个性质可以扩展到无穷个集合的交：如果对于任意 $\alpha\in\mathcal{A}$，$S_\alpha$ 都是凸的，那么 $\bigcap_{\alpha\in\mathcal{A}}S_\alpha$ 也是凸集。（子空间、仿射集合和凸锥对于任意交运算也是封闭的。）作为一个简单的例子，多面体是半空间和超平面（它们都是凸集）的交集，因而是凸的。

\noindent\textbf{例 2.7}\quad 半正定锥 $\mathbf{S}_+^n$ 可以表示为，
\[
  \bigcap_{z\ne 0}\{X\in\mathbf{S}^n\mid z^TXz\ge 0\}.
\]
对于任意 $z\ne 0$，$z^TXz$ 是关于 $X$ 的（不恒等于零的）线性函数，因此集合
\[
  \{X\in\mathbf{S}^n\mid z^TXz\ge 0\}
\]
实际上就是 $\mathbf{S}^n$ 的半空间。由此可见，半正定锥是无穷个半空间的交集，因此是凸的。

\noindent\textbf{例 2.8}\quad 考虑集合
\begin{equation}
  S=\{x\in\mathbb{R}^m\mid |p(t)|\le 1\text{ 对于 }|t|\le \pi/3\},
  \tag{2.10}
\end{equation}
其中 $p(t)=\sum_{k=1}^m x_k\cos kt$。集合 $S$ 可以表示为无穷个平板的交集：$S=\bigcap_{|t|\le\pi/3}S_t$，其中，
\[
  S_t=\{x\mid -1\le (\cos t,\cdots,\cos mt)^Tx\le 1\},
\]
因此，$S$ 是凸的。对于 $m=2$ 的情况，它的定义和集合可见图2.13和图2.14。

在上面这些例子中，我们通过将集合表示为（可能无穷多个）半空间的交集来表明集合的凸性。反过来，在 §2.5.1 中，我们也将看到：每一个闭的凸集 $S$ 是（通常为无限多个）半空间的交集。事实上，一个闭集 $S$ 是包含它的所有半空间的交集：
\[
  S=\bigcap\{\mathcal{H}\mid \mathcal{H}\text{ 是半空间},\ S\subseteq\mathcal{H}\}.
\]

\subsection*{2.3.2\quad 仿射函数}

函数 $f:\mathbb{R}^n\to\mathbb{R}^m$ 是仿射的，如果它是一个线性函数和一个常数的和，即具有 $f(x)=Ax+b$ 的形式，其中 $A\in\mathbb{R}^{m\times n}$，$b\in\mathbb{R}^m$。假设 $S\subseteq\mathbb{R}^n$ 是凸的，并且 $f:\mathbb{R}^n\to\mathbb{R}^m$ 是仿射函数。那么，$S$ 在 $f$ 下的象
\[
  f(S)=\{f(x)\mid x\in S\}
\]
是凸的。类似地，如果 $f:\mathbb{R}^k\to\mathbb{R}^n$ 是仿射函数，那么 $S$ 在 $f$ 下的原象
\[
  f^{-1}(S)=\{x\mid f(x)\in S\}
\]
是凸的。

两个简单的例子是伸缩和平移。如果 $S\subseteq\mathbb{R}^n$ 是凸集，$\alpha\in\mathbb{R}$ 并且 $a\in\mathbb{R}^n$，那么，集合 $\alpha S$ 和 $S+a$ 是凸的，其中
\[
  \alpha S=\{\alpha x\mid x\in S\},\qquad S+a=\{x+a\mid x\in S\}.
\]
一个凸集向它的某几个坐标的投影是凸的，即：如果 $S\subseteq\mathbb{R}^m\times\mathbb{R}^n$ 是凸集，那么
\[
  T=\{x_1\in\mathbb{R}^m\mid (x_1,x_2)\in S\text{ 对于某些 }x_2\in\mathbb{R}^n\}
\]
是凸集。

两个集合的和可以定义为：
\[
  S_1+S_2=\{x+y\mid x\in S_1,\ y\in S_2\}.
\]
如果 $S_1$ 和 $S_2$ 是凸集，那么，$S_1+S_2$ 是凸的。可以看出，如果 $S_1$ 和 $S_2$ 是凸的，那么其直积或 Cartesian 乘积
\[
  S_1\times S_2=\{(x_1,x_2)\mid x_1\in S_1,\ x_2\in S_2\}
\]
也是凸集。这个集合在线性函数 $f(x_1,x_2)=x_1+x_2$ 下的象是和 $S_1+S_2$。

我们也可以考虑 $S_1,S_2\in\mathbb{R}^n\times\mathbb{R}^m$ 的部分和，定义为
\[
  S=\{(x,y_1+y_2)\mid (x,y_1)\in S_1,\ (x,y_2)\in S_2\},
\]
其中 $x\in\mathbb{R}^n,\ y_i\in\mathbb{R}^m$。$m=0$ 时，部分和给出了 $S_1$ 和 $S_2$ 的交集；$n=0$ 时，部分和等于集合之和。凸集的部分和是凸集（参见习题 2.16）。

\noindent\textbf{例 2.9}\quad 多面体。多面体 $\{x\mid Ax\preceq b,\ Cx=d\}$ 可以表示为非负象限和原点的 Cartesian 乘积在仿射函数 $f(x)=(b-Ax,d-Cx)$ 下的原象：
\[
  \{x\mid Ax\preceq b,\ Cx=d\}
  =\{x\mid f(x)\in\mathbb{R}_+^m\times\{0\}\}.
\]

\noindent\textbf{例 2.10}\quad 线性矩阵不等式的解。条件：
\begin{equation}
  A(x)=x_1A_1+\cdots+x_nA_n\preceq B
  \tag{2.11}
\end{equation}
称为关于 $x$ 的线性矩阵不等式（LMI），其中 $B,A_i\in\mathbf{S}^m$。（注意它与有序线性不等式
\[
  a^Tx=x_1a_1+\cdots+x_na_n\le b
\]
的相似性，其中 $b,a_i\in\mathbb{R}$。）

线性矩阵不等式的解 $\{x\mid A(x)\preceq B\}$ 是凸集。事实上，它是半正定锥在由 $f(x)=B-A(x)$ 给定的仿射映射 $f:\mathbb{R}^n\to\mathbf{S}^m$ 下的原象。

\noindent\textbf{例 2.11}\quad 双曲锥。集合
\[
  \{x\mid x^TPx\le (c^Tx)^2,\ c^Tx\ge 0\}
\]
是凸集，其中 $P\in\mathbf{S}_+^n,\ c\in\mathbb{R}^n$。这是因为它是二阶锥
\[
  \{(z,t)\mid z^Tz\le t^2,\ t\ge 0\}
\]
在仿射函数 $f(x)=(P^{1/2}x,c^Tx)$ 下的原象。

\noindent\textbf{例 2.12}\quad 椭球。椭球
\[
  \mathcal{E}=\{x\mid (x-x_c)^TP^{-1}(x-x_c)\le 1\}
\]
是单位 Euclid 球 $\{u\mid \|u\|_2\le 1\}$ 在仿射映射 $f(u)=P^{1/2}u+x_c$ 下的象，其中 $P\in\mathbf{S}_{++}^n$。（同时也是单位球在仿射映射 $g(x)=P^{-1/2}(x-x_c)$ 下的原象。）

\subsection*{2.3.3\quad 线性分式及透视函数}

本节将讨论一类称为线性分式的函数，它比仿射函数更普遍，并且仍然保凸。

\noindent\textbf{透视函数}

我们定义 $P:\mathbb{R}^{n+1}\to\mathbb{R}^n$，$P(z,t)=z/t$ 为透视函数，其定义域为 $\mathbf{dom}\,P=\mathbb{R}^n\times\mathbb{R}_{++}$。（此处的 $\mathbb{R}_{++}$ 表示正实数集合，即 $\mathbb{R}_{++}=\{x\in\mathbb{R}\mid x>0\}$。）透视函数对向量进行伸缩，或称为规范化，使得最后一维分量为 1 并舍弃之。

\noindent\textbf{注释 2.1}\quad 我们用小孔成象来解释透视函数。（$\mathbb{R}^3$ 中的）小孔照相机由一个不透明的水平面 $x_3=0$ 和一个在原点的小孔组成，光线通过这个小孔在 $x_3=-1$ 呈现出一个水平图像。在相机上方 $x$（$x_3>0$）处的一个物体，在相平面的点 $-(x_1/x_3,x_2/x_3,1)$ 处形成一个图像。忽略象点的最后一维分量（因为它恒等于 $-1$），$x$ 处的点的象在象平面上呈现于 $y=-(x_1/x_3,x_2/x_3)=-P(x)$ 处。图2.15显示了这个过程。

如果 $C\subseteq\mathbf{dom}\,P$ 是凸集，那么它的象
\[
  P(C)=\{P(x)\mid x\in C\}
\]
也是凸集。这个结论很直观：通过小孔观察一个凸的物体，可以得到凸的象。为解释这个事实，我们将说明在透视函数作用下，线段将被映射成线段。（也可以这样理解，通过小孔，一条线段的象是一条线段。）假设 $x=(\tilde{x},x_{n+1})$，$y=(\tilde{y},y_{n+1})\in\mathbb{R}^{n+1}$，并且 $x_{n+1}>0,\ y_{n+1}>0$。那么，对于 $0\le\theta\le 1$，
\[
  P(\theta x+(1-\theta)y)
  =
  \frac{\theta\tilde{x}+(1-\theta)\tilde{y}}{\theta x_{n+1}+(1-\theta)y_{n+1}}
  =\mu P(x)+(1-\mu)P(y),
\]
其中，
\[
  \mu=\frac{\theta x_{n+1}}{\theta x_{n+1}+(1-\theta)y_{n+1}}\in[0,1].
\]
$\theta$ 和 $\mu$ 之间的关系是单调的：当 $\theta$ 在 0、1 间变化时（形成线段 $[x,y]$），$\mu$ 也在 0、1 间变化（形成线段 $[P(x),P(y)]$）。这说明 $P([x,y])=[P(x),P(y)]$。

现在假设 $C$ 是凸的，并且有 $C\subseteq\mathbf{dom}\,P$（即对于所有 $x\in C,\ x_{n+1}>0$）及 $x,y\in C$。为显示 $P(C)$ 的凸性，我们需要说明线段 $[P(x),P(y)]$ 在 $P(C)$ 中。这条线段是线段 $[x,y]$ 在 $P$ 的象，因而属于 $P(C)$。

一个凸集在透视函数下的原象也是凸的：如果 $C\subseteq\mathbb{R}^n$ 为凸集，那么
\[
  P^{-1}(C)=\{(x,t)\in\mathbb{R}^{n+1}\mid x/t\in C,\ t>0\}
\]
是凸集。为证明这点，假设 $(x,t)\in P^{-1}(C)$，$(y,s)\in P^{-1}(C)$，$0\le\theta\le 1$。我们需要说明
\[
  \theta(x,t)+(1-\theta)(y,s)\in P^{-1}(C),
\]
即
\[
  \frac{\theta x+(1-\theta)y}{\theta t+(1-\theta)s}\in C
\]
（显然地，$\theta t+(1-\theta)s>0$）。这可从下式看出，
\[
  \frac{\theta x+(1-\theta)y}{\theta t+(1-\theta)s}
  =\mu(x/t)+(1-\mu)(y/s),
\]
其中，
\[
  \mu=\frac{\theta t}{\theta t+(1-\theta)s}\in[0,1].
\]

\noindent\textbf{线性分式函数}

线性分式函数由透视函数和仿射函数复合而成。设 $g:\mathbb{R}^n\to\mathbb{R}^{m+1}$ 是仿射的，即
\begin{equation}
  g(x)=
  \begin{bmatrix}
    A\\
    c^T
  \end{bmatrix}x+
  \begin{bmatrix}
    b\\
    d
  \end{bmatrix},
  \tag{2.12}
\end{equation}
其中 $A\in\mathbb{R}^{m\times n}$，$b\in\mathbb{R}^m$，$c\in\mathbb{R}^n$ 并且 $d\in\mathbb{R}$。则由 $f=P\circ g$ 给出的函数 $f:\mathbb{R}^n\to\mathbb{R}^m$
\begin{equation}
  f(x)=(Ax+b)/(c^Tx+d),\qquad
  \mathbf{dom}\,f=\{x\mid c^Tx+d>0\},
  \tag{2.13}
\end{equation}
称为线性分式（或投射）函数。如果 $c=0,\ d>0$，则 $f$ 的定义域为 $\mathbb{R}^n$，并且 $f$ 是仿射函数。因此，我们可以将仿射和线性函数视为特殊的线性分式函数。

\noindent\textbf{注释 2.2}\quad 投射解释。可以很方便地将一个线性分式函数表示为：将矩阵
\begin{equation}
  Q=
  \begin{bmatrix}
    A & b\\
    c^T & d
  \end{bmatrix}\in\mathbb{R}^{(m+1)\times(n+1)}
  \tag{2.14}
\end{equation}
作用于（即左乘）点 $(x,1)$，得到 $(Ax+b,c^Tx+d)$；然后将所得结果做伸缩变换或归一化，以使得其最后一个分量为一，得到 $(f(x),1)$。

也可以从几何上进行解释，即这个表达式将 $\mathbb{R}^n$ 与 $\mathbb{R}^{n+1}$ 空间上的一组射线联系了起来。也就是说，对于 $\mathbb{R}^n$ 空间的一个点 $z$，我们可以构造一个 $\mathbb{R}^{n+1}$ 空间的（开）射线 $\mathcal{P}(z)=\{t(z,1)\mid t>0\}$。这条射线上每个点的最后一个分量均是正值。反之，$\mathbb{R}^{n+1}$ 空间中每条以原点为顶点并且最后一个分量为正值的射线均可以由一些 $v\in\mathbb{R}^n$ 表示为 $\mathcal{P}(v)=\{t(v,1)\mid t\ge 0\}$。$\mathcal{P}$ 表示了 $\mathbb{R}^n$ 与最后一个分量为正的射线之间的（投射）关系，这种关系是一一对应和满的。

线性分式函数 (2.13) 可以表示为：
\[
  f(x)=\mathcal{P}^{-1}(Q\mathcal{P}(x)).
\]
因此，从 $x\in\mathbf{dom}\,f$ 出发，即由一个满足 $c^Tx+d>0$ 的点，我们可以得到 $\mathbb{R}^{n+1}$ 空间的一条射线 $\mathcal{P}(x)$。将线性变换矩阵 $Q$ 作用于这条射线，就可以得到另一条射线 $Q\mathcal{P}(x)$。因为 $x\in\mathbf{dom}\,f$，这条射线的最后一个分量为正。最后，我们也可以通过逆投射变换恢复出 $f(x)$。

类似于透视函数，线性分式函数也是保凸的。如果 $C$ 是凸集并且在 $f$ 的定义域中（即任意 $x\in C$ 满足 $c^Tx+d>0$），那么 $C$ 的象 $f(C)$ 也是凸集。根据前述的结果可以直接得到这个结论：$C$ 在仿射映射 (2.12) 下的象是凸的，并且在透视函数 $P$ 下的映射（即 $f(C)$）是凸的。类似地，如果 $C\subseteq\mathbb{R}^m$ 是凸集，那么其原象 $f^{-1}(C)$ 也是凸的。

\noindent\textbf{例 2.13}\quad 条件概率。设 $u$ 和 $v$ 是分别在 $\{1,\cdots,n\}$ 和 $\{1,\cdots,m\}$ 中取值的随机变量，并且 $p_{ij}$ 表示概率 $\operatorname{prob}(u=i,v=j)$。那么条件概率 $f_{ij}=\operatorname{prob}(u=i\mid v=j)$ 由下式给出
\[
  f_{ij}=\frac{p_{ij}}{\sum_{k=1}^n p_{kj}}.
\]
因此，$f$ 可以通过一个线性分式映射从 $p$ 得到。

可以知道，如果 $C$ 是一个关于 $(u,v)$ 的联合密度的凸集，那么相应的 $u$ 的条件密度（给定 $v$）的集合也是凸集。

图2.16表示了集合 $C\subseteq\mathbb{R}^2$ 及其在下面的线性分式函数下的象：
\[
  f(x)=\frac{1}{x_1+x_2+1}x,\qquad
  \mathbf{dom}\,f=\{(x_1,x_2)\mid x_1+x_2+1>0\}.
\]

\section*{2.4\quad 广义不等式}

\subsection*{2.4.1\quad 正常锥与广义不等式}

称锥 $K\subseteq\mathbb{R}^n$ 为正常锥，如果它满足下列条件
\begin{itemize}
  \item $K$ 是凸的。
  \item $K$ 是闭的。
  \item $K$ 是实的，即具有非空内部。
  \item $K$ 是尖的，即不包含直线（或者等价地，$x\in K,\ -x\in K\Longrightarrow x=0$）。
\end{itemize}
正常锥 $K$ 可以用来定义广义不等式，即 $\mathbb{R}^n$ 上的偏序关系。这种偏序关系和 $\mathbb{R}$ 上的标准序有很多相同的性质。用正常锥 $K$ 可以定义 $\mathbb{R}^n$ 上的偏序关系如下
\[
  x\preceq_K y\Longleftrightarrow y-x\in K.
\]
$y\succeq_K x$ 也可以写为 $x\preceq_K y$。类似地，我们定义相应的严格偏序关系为
\[
  x\prec_K y\Longleftrightarrow y-x\in\mathbf{int}\,K,
\]
并且可以同样地定义 $x\succ_K y$。（为将广义不等式 $\preceq_K$ 与严格的广义不等式区分开，我们有时也称 $\preceq_K$ 为不严格的广义不等式）。

当 $K=\mathbb{R}_+$ 时，偏序关系 $\preceq_K$ 就是通常意义上 $\mathbb{R}$ 中的序 $\le$；相应地，严格偏序关系 $\prec_K$ 与 $\mathbb{R}$ 上的严格序 $<$ 相同。因此，广义不等式包含了 $\mathbb{R}$ 上的（不严格和严格）不等式，它是广义不等式的一种特殊情况。

\noindent\textbf{例 2.14}\quad 非负象限及分量不等式。非负象限 $K=\mathbb{R}_+^n$ 是一个正常锥。相应的广义不等式 $\preceq_K$ 对应于向量间的分量不等式，即 $x\preceq_K y$ 等价于 $x_i\le y_i,\ i=1,\cdots,n$。相应地，其严格不等式对应于严格的分量不等式，即 $x\prec_K y$ 等价于 $x_i<y_i,\ i=1,\cdots,n$。

我们将常常使用对应于非负象限的不严格和严格的偏序关系，因此省略下标 $\mathbb{R}_+^n$。当 $\preceq$ 或 $\prec$ 出现在向量间的时候，该符号应被理解为分量不等式。

\noindent\textbf{例 2.15}\quad 半正定锥和矩阵不等式。半正定锥是 $\mathbf{S}^n$ 空间中的正常锥，相应的广义不等式 $\preceq_K$ 就是通常的矩阵不等式，即 $X\preceq_KY$ 等价于 $Y-X$ 为半正定矩阵。（在 $\mathbf{S}^n$ 中）$\mathbf{S}_+^n$ 的内部由正定矩阵组成，因此严格广义不等式也等同于通常的对称矩阵的严格不等式，即 $X\prec_KY$ 等价于 $Y-X$ 为正定矩阵。

这里，也是由于经常使用这种偏序关系，因此省略其下标，即对于对称矩阵，我们将广义不等式简写为 $X\preceq Y$ 或 $X\prec Y$，它们表示关于半正定锥的广义不等式。

\noindent\textbf{例 2.16}\quad $[0,1]$ 上非负的多项式锥。$K$ 定义如下
\begin{equation}
  K=\{c\in\mathbb{R}^n\mid c_1+c_2t+\cdots+c_nt^{n-1}\ge 0\text{ 对于 }t\in[0,1]\},
  \tag{2.15}
\end{equation}
即 $K$ 是 $[0,1]$ 上最高 $n-1$ 阶的非负多项式（系数）锥。可以看出 $K$ 是一个正常锥，其内部是 $[0,1]$ 上为正的多项式的系数集合。

两个向量 $c,d\in\mathbb{R}^n$ 满足 $c\preceq_K d$ 的充要条件是，对于所有 $t\in[0,1]$ 有
\[
  c_1+c_2t+\cdots+c_nt^{n-1}
  \le d_1+d_2t+\cdots+d_nt^{n-1}.
\]

\noindent\textbf{广义不等式的性质}

广义不等式 $\preceq_K$ 满足许多性质，例如
\begin{itemize}
  \item $\preceq_K$ 对于加法是保序的：如果 $x\preceq_K y$ 并且 $u\preceq_K v$，那么 $x+u\preceq_K y+v$。
  \item $\preceq_K$ 具有传递性：如果 $x\preceq_K y$ 并且 $y\preceq_K z$，那么 $x\preceq_K z$。
  \item $\preceq_K$ 对于非负数乘是保序的：如果 $x\preceq_K y$ 并且 $\alpha\ge 0$，那么 $\alpha x\preceq_K\alpha y$。
  \item $\preceq_K$ 是自反的：$x\preceq_K x$。
  \item $\preceq_K$ 是反对称的：如果 $x\preceq_K y$ 并且 $y\preceq_K x$，那么 $x=y$。
  \item $\preceq_K$ 对于极限运算是保序的：如果对于 $i=1,2,\cdots$ 均有 $x_i\preceq_K y_i$，当 $i\to\infty$ 时，有 $x_i\to x$ 和 $y_i\to y$，那么 $x\preceq_K y$。
\end{itemize}
相应的广义不等式 $\prec_K$ 也满足一些性质，例如
\begin{itemize}
  \item 如果 $x\prec_K y$，那么 $x\preceq_K y$。
  \item 如果 $x\prec_K y$ 并且 $u\preceq_K v$，那么 $x+u\prec_K y+v$。
  \item 如果 $x\prec_K y$ 并且 $\alpha>0$，那么 $\alpha x\prec_K\alpha y$。
  \item $x\not\prec_K x$。
  \item 如果 $x\prec_K y$，那么对于足够小的 $u$ 和 $v$ 有 $x+u\prec_K y+v$。
\end{itemize}
这些性质可以从 $\preceq_K$ 和 $\prec_K$ 的定义以及正常锥的性质中直接得到，参见习题 2.30。

\subsection*{2.4.2\quad 最小与极小元}

广义不等式的符号 $(\preceq_K,\prec_K)$ 似乎表明它们与 $\mathbb{R}$ 上的普通不等式 $(\le,<)$ 有着相同的性质。虽然普通不等式的许多性质对于广义不等式确实成立，但很多重要的性质并不如此。最明显的区别在于，$\mathbb{R}$ 上的 $\le$ 是一个线性序，即任意两点都是可比的，也就是说 $x\le y$ 和 $y\le x$ 二者必居其一。这个性质对于其他广义不等式并不成立。这导致了最小、最大这些概念在广义不等式环境下变得更加复杂。本节将对此进行简要的讨论。

如果对于每个 $y\in S$，均有 $x\preceq_K y$，我们称 $x\in S$ 是 $S$（关于广义不等式 $\preceq_K$）的最小元。类似地，我们可以定义关于广义不等式的最大元。如果一个集合有最小（或最大）元，那么它们是唯一的。相对应的概念是极小元。如果 $y\in S,\ y\preceq_K x$ 可以推得 $y=x$，那么我们称 $x\in S$ 是 $S$ 上（关于广义不等式 $\preceq_K$）的极小元。同样地，可以定义极大元。一个集合可以有多个极小（或极大）元。

用简单的集合符号，我们可以对最小元和极小元进行描述。元素 $x\in S$ 是 $S$ 中的一个最小元，当且仅当
\[
  S\subseteq x+K.
\]
这里 $x+K$ 表示可以与 $x$ 相比并且大于或等于（根据 $\preceq_K$）$x$ 的所有元素。元素 $x\in S$ 是极小元，当且仅当
\[
  (x-K)\cap S=\{x\}.
\]
这里 $x-K$ 表示可以与 $x$ 相比并且小于或等于（根据 $\preceq_K$）$x$ 的所有元素，它与 $S$ 的唯一共同点即是 $x$。

$K=\mathbb{R}_+$ 导出的实际上就是 $\mathbb{R}$ 上一般的序。此时，极小和最小的概念是一致的，也符合集合最小元素的通常定义。

\noindent\textbf{例 2.17}\quad 考虑锥 $\mathbb{R}_+^2$，它导出的是 $\mathbb{R}^2$ 上的关于分量的不等式。对此，我们可以给出一些关于极小元和最小元的简单的几何描述。不等式 $x\preceq y$ 的含义是 $y$ 在 $x$ 之上、之右。$x\in S$ 是集合 $S$ 的最小元，表明 $S$ 的其他所有点都在它之上、之右。而 $x$ 为集合 $S$ 的极小元，是指 $S$ 中没有任何一个点在 $x$ 之下、之左，其区别可见图2.17。

\noindent\textbf{例 2.18}\quad 对称矩阵集合中的最小元和极小元。我们用 $A\in\mathbf{S}_{++}^n$ 表示一个圆心在原点的椭圆，即
\[
  \mathcal{E}_A=\{x\mid x^TA^{-1}x\le 1\}.
\]
我们知道 $A\preceq B$ 等价于 $\mathcal{E}_A\subseteq\mathcal{E}_B$。

给定 $v_1,\cdots,v_k\in\mathbb{R}^n$ 并定义
\[
  S=\{P\in\mathbf{S}_{++}^n\mid v_i^TP^{-1}v_i\le 1,\ i=1,\cdots,k\}.
\]
它对应于包含了点 $v_1,\cdots,v_k$ 的椭圆的集合。集合 $S$ 没有最小元：对于任意包含点 $v_1,\cdots,v_k$ 的椭圆，我们总可以找到另一个包含这些点但不可比的椭圆。一个椭圆是极小的，如果它包含这些点但没有更小的椭圆也包含这些点。图2.18显示了 $\mathbb{R}^2$ 上 $k=2$ 时的一个例子。

\section*{2.5\quad 分离与支撑超平面}

\subsection*{2.5.1\quad 超平面分离定理}

本节中我们将阐述一个在之后非常重要的想法：用超平面或仿射函数将两个不相交的凸集分离开来。其基本结果就是超平面分离定理：假设 $C$ 和 $D$ 是两个不相交的凸集，即 $C\cap D=\emptyset$，那么存在 $a\ne 0$ 和 $b$ 使得对于所有 $x\in C$ 有 $a^Tx\le b$，对于所有 $x\in D$ 有 $a^Tx\ge b$。换言之，仿射函数 $a^Tx-b$ 在 $C$ 中非正，而在 $D$ 中非负。超平面 $\{x\mid a^Tx=b\}$ 称为集合 $C$ 和 $D$ 的分离超平面，或称超平面分离了集合 $C$ 和 $D$，参见图2.19。

\noindent\textbf{超平面分离定理的证明}

这里我们考虑一个特殊的情况，并将一般情况的证明留作扩展习题（习题 2.22）。我们假设 $C$ 和 $D$ 的（Euclid）距离为正，这里的距离定义为
\[
  \mathbf{dist}(C,D)=\inf\{\|u-v\|_2\mid u\in C,\ v\in D\},
\]
并且存在 $c\in C$ 和 $d\in D$ 达到这个最小距离，即 $\|c-d\|_2=\mathbf{dist}(C,D)$。（这些条件是可以被满足的，例如当 $C$ 和 $D$ 是闭的并且其中之一是有界的。）

定义
\[
  a=d-c,\qquad b=\frac{\|d\|_2^2-\|c\|_2^2}{2}.
\]
我们将显示仿射函数
\[
  f(x)=a^Tx-b=(d-c)^T(x-(1/2)(d+c))
\]
在 $C$ 中非正而在 $D$ 中非负，即超平面 $\{x\mid a^Tx=b\}$ 分离了 $C$ 和 $D$。这个超平面与连接 $c$ 和 $d$ 之间的线段相垂直并且穿过其中点，如图2.20所示。

我们首先证明 $f$ 在 $D$ 中非负。关于 $f$ 在 $C$ 中非正的证明是相似的（只需将 $C$ 和 $D$ 交换并考虑 $-f$ 即可）。假设存在一个点 $u\in D$，并且
\begin{equation}
  f(u)=(d-c)^T(u-(1/2)(d+c))<0.
  \tag{2.16}
\end{equation}
我们可以将 $f(u)$ 表示为
\[
  f(u)=(d-c)^T(u-d+(1/2)(d-c))=(d-c)^T(u-d)+(1/2)\|d-c\|_2^2.
\]
可以看出式 (2.16) 意味着 $(d-c)^T(u-d)<0$。于是，我们观察到
\[
  \left.\frac{d}{dt}\|d+t(u-d)-c\|_2^2\right|_{t=0}
  =2(d-c)^T(u-d)<0.
\]
因此，对于足够小的 $t>0$ 及 $t\le 1$，我们有
\[
  \|d+t(u-d)-c\|_2<\|d-c\|_2,
\]
即点 $d+t(u-d)$ 比 $d$ 更靠近 $c$。因为 $D$ 是包含 $d$ 和 $u$ 的凸集，我们有 $d+t(u-d)\in D$。但这是不可能的，因为根据假设，$d$ 应当是 $D$ 中离 $C$ 最近的点。

\noindent\textbf{例 2.19}\quad 仿射集与凸集的分离。设 $C$ 是凸集，而 $D$ 是仿射的，即 $D=\{Fu+g\mid u\in\mathbb{R}^m\}$，其中 $F\in\mathbb{R}^{n\times m}$。设 $C$ 和 $D$ 不相交，那么根据超平面分离定理，存在 $a\ne 0$ 和 $b$ 使得对于所有 $x\in C$ 有 $a^Tx\le b$，对于所有 $x\in D$ 有 $a^Tx\ge b$。

这里 $a^Tx\ge b$ 对于所有 $x\in D$ 均成立，表明对于任意 $u\in\mathbb{R}^m$ 均有 $a^TFu\ge b-a^Tg$。但是在 $\mathbb{R}^m$ 上，只有当一个线性函数为零时，它才是有界的。因此，可以推知 $a^TF=0$（并且因此有 $b\le a^Tg$）。

所以，我们可知存在 $a\ne 0$ 使得 $F^Ta=0$ 和 $a^Tx\le a^Tg$ 对于所有 $x\in C$ 均成立。

\noindent\textbf{严格分离}

如果之前构造的分离超平面满足更强的条件，即对于任意 $x\in C$ 有 $a^Tx<b$ 并且对于任意 $x\in D$ 有 $a^Tx>b$，则称其为集合 $C$ 和 $D$ 的严格分离。简单的例子就可以看出，对于一般的情况，不相交的凸集并不一定能够被超平面严格分离（即使集合是闭集，参见习题 2.23）。但是，在很多特殊的情况下，可以构造严格分离。

\noindent\textbf{例 2.20}\quad 点和闭凸集的严格分离。令 $C$ 为闭凸集，而 $x_0\notin C$，那么存在将 $x_0$ 与 $C$ 严格分离的超平面。

为说明这一点，需要注意，对于足够小的 $\epsilon>0$，存在两个不相交的集合 $C$ 和 $B(x_0,\epsilon)$。根据超平面分离定理，存在 $a\ne0$ 和 $b$，使得对于任意 $x\in C$ 有 $a^Tx\le b$；对于任意 $x\in B(x_0,\epsilon)$ 有 $a^Tx\ge b$。

利用 $B(x_0,\epsilon)=\{x_0+u\mid \|u\|_2\le\epsilon\}$，可以将前述第二个条件表示为
\[
  a^T(x_0+u)\ge b\quad \text{对于所有 }\|u\|_2\le\epsilon.
\]
$u=-\epsilon a/\|a\|_2$ 极小化了上式的左端，代入可得
\[
  a^Tx_0-\epsilon\|a\|_2\ge b.
\]
所以，仿射函数
\[
  f(x)=a^Tx-b-\epsilon\|a\|_2/2
\]
在 $C$ 上是负的，而在 $x_0$ 点是正的。

作为一个直接的结果，我们可以得到前面已提及的事实：一个闭凸集是包含它的所有半空间的交集。事实上，令 $C$ 为闭和凸的，$S$ 为所有包含 $C$ 的半空间。显然，$x\in C\Rightarrow x\in S$。为证明反方向，假设存在 $x\in S$ 并且 $x\notin C$。根据严格分离的结果，存在一个将 $x$ 与 $C$ 严格分离的超平面，即存在一个包含 $C$ 但不包含 $x$ 的半空间。也就是说，$x\notin S$。

\noindent\textbf{超平面分离定理的逆定理}

超平面分离定理的逆定理（即分离超平面的存在表明 $C$ 和 $D$ 不相交）是不成立的，除非在凸性之外再给 $C$ 或 $D$ 附加其他约束。作为一个简单的反例，我们考虑 $C=D=\{0\}\subseteq\mathbb{R}$，超平面 $x=0$ 可以分离 $C$ 和 $D$。

通过给 $C$ 和 $D$ 增加一些条件，可以得到超平面分离定理的多种逆定理。作为一个简单的例子，设 $C$ 和 $D$ 是凸集，$C$ 是开集，如果存在一个仿射函数 $f$，它在 $C$ 中非正而在 $D$ 中非负，那么 $C$ 和 $D$ 不相交。（为说明此结论，首先可知 $f$ 在 $C$ 上是负的。否则，如果 $f$ 在 $C$ 中的某一点为零，那么 $f$ 在这个点附近会取得正值，这与前述矛盾。因此 $C$ 和 $D$ 一定是不相交的，因为 $f$ 在 $C$ 中为负，而在 $D$ 中非负。）将逆定理与超平面分离定理相结合，我们可以得到下面的结论：任何两个凸集 $C$ 和 $D$，如果其中至少有一个是开集，那么当且仅当存在分离超平面时，它们不相交。

\noindent\textbf{例 2.21}\quad 严格线性不等式的择一定理。我们导出严格线性不等式
\begin{equation}
  Ax\prec b
  \tag{2.17}
\end{equation}
有解的充要条件。该不等式不可行的充要条件是（凸）集
\[
  C=\{b-Ax\mid x\in\mathbb{R}^n\},\qquad
  D=\mathbb{R}_{++}^m=\{y\in\mathbb{R}^m\mid y\succ0\}
\]
不相交。集合 $D$ 是开集，而 $C$ 是仿射集合。根据前述的结论，$C$ 和 $D$ 不相交的充要条件是，存在分离超平面，即存在非零的 $\lambda\in\mathbb{R}^m$ 和 $\mu\in\mathbb{R}$ 使得 $C$ 中 $\lambda^Ty\le\mu$ 而 $D$ 中 $\lambda^Ty\ge\mu$。

这些条件可以被简化。第一个条件意味着对于所有 $x$ 都有 $\lambda^T(b-Ax)\le\mu$。这表明（如例2.19所示）$A^T\lambda=0$，$\lambda^Tb\le\mu$。第二个不等式意味着 $\lambda^Ty\ge\mu$ 对于所有 $y\succ0$ 均成立。这表明 $\mu\le0$ 且 $\lambda\succeq0,\lambda\ne0$。

将这些结果放在一起，我们可以得知严格不等式组 (2.17) 无解的充要条件是存在 $\lambda\in\mathbb{R}^m$ 使得
\begin{equation}
  \lambda\ne0,\qquad \lambda\succeq0,\qquad A^T\lambda=0,\qquad \lambda^Tb\le0.
  \tag{2.18}
\end{equation}
这些不等式和等式关于 $\lambda\in\mathbb{R}^m$ 也是线性的。我们称式(2.17) 和式(2.18) 构成一对择一选择：对于任意的 $A$ 和 $b$，两者中仅有一组有解。

\subsection*{2.5.2\quad 支撑超平面}

设 $C\subseteq\mathbb{R}^n$ 而 $x_0$ 是其边界 $\mathbf{bd}\,C$ 上的一点，即
\[
  x_0\in\mathbf{bd}\,C=\mathbf{cl}\,C\setminus\mathbf{int}\,C.
\]
如果 $a\ne0$，并且对任意 $x\in C$ 满足 $a^Tx\le a^Tx_0$，那么称超平面 $\{x\mid a^Tx=a^Tx_0\}$ 为集合 $C$ 在点 $x_0$ 处的支撑超平面。这等于说点 $x_0$ 与集合 $C$ 被超平面所分离 $\{x\mid a^Tx=a^Tx_0\}$。其几何解释是超平面 $\{x\mid a^Tx=a^Tx_0\}$ 与 $C$ 相切于点 $x_0$，而且半空间 $\{x\mid a^Tx\le a^Tx_0\}$ 包含 $C$，参见图2.21。

一个基本的结论，称为支撑超平面定理，表明对于任意非空的凸集 $C$ 和任意 $x_0\in\mathbf{bd}\,C$，在 $x_0$ 处存在 $C$ 的支撑超平面。支撑超平面定理从超平面分离定理很容易得到证明。需要区分两种情况。如果 $C$ 的内部非空，对于 $\{x_0\}$ 和 $\mathbf{int}\,C$ 应用超平面分离定理可以直接得到所需的结论。如果 $C$ 的内部是空集，则 $C$ 必处于小于 $n$ 维的一个仿射集合中，并且任意包含这个仿射集合的超平面一定包含 $C$ 和 $x_0$，这是一个（平凡的）支撑超平面。

支撑超平面定理也有一个不完全的逆定理：如果一个集合是闭的，具有非空内部，并且其边界上每个点均存在支撑超平面，那么它是凸的（参见习题 2.27）。

\section*{2.6\quad 对偶锥与广义不等式}

\subsection*{2.6.1\quad 对偶锥}

令 $K$ 为一个锥。集合
\begin{equation}
  K^*=\{y\mid x^Ty\ge0,\ \forall x\in K\}
  \tag{2.19}
\end{equation}
称为 $K$ 的对偶锥。顾名思义，$K^*$ 是一个锥，并且它总是凸的，即使 $K$ 不是凸锥（参见习题 2.31）。

从几何上看，$y\in K^*$ 当且仅当 $-y$ 是 $K$ 在原点的一个支撑超平面的法线，如图2.22所示。

\noindent\textbf{例 2.22}\quad 子空间。子空间 $V\subseteq\mathbb{R}^n$（这是一个锥）的对偶锥是其正交补 $V^\perp=\{y\mid y^Tv=0,\ \forall v\in V\}$。

\noindent\textbf{例 2.23}\quad 非负象限。锥 $\mathbb{R}_+^n$ 的对偶是它本身：
\[
  y^Tx\ge0,\ \forall x\succeq0\Longleftrightarrow y\succeq0.
\]
我们称这种锥自对偶。

\noindent\textbf{例 2.24}\quad 半正定锥。在 $n\times n$ 对称矩阵的集合 $\mathbf{S}^n$ 上，我们使用其标准内积 $\operatorname{tr}(XY)=\sum_{i,j=1}^n X_{ij}Y_{ij}$（参见附录 A.1.1）。半正定锥 $\mathbf{S}_+^n$ 是自对偶的，即对于任意的 $X,Y\in\mathbf{S}^n$，
\[
  \operatorname{tr}(XY)\ge0,\ \forall X\succeq0 \Longleftrightarrow Y\succeq0.
\]
下面，我们将说明这一结论。

假设 $Y\notin\mathbf{S}_+^n$，那么存在 $q\in\mathbb{R}^n$ 并且
\[
  q^TYq=\operatorname{tr}(qq^TY)<0.
\]
于是半正定矩阵 $X=qq^T$ 满足 $\operatorname{tr}(XY)<0$，由此可知 $Y\notin(\mathbf{S}_+^n)^*$。

假设 $X,Y\in\mathbf{S}_+^n$，我们可以利用特征值分解将 $X$ 表述为 $X=\sum_{i=1}^n\lambda_iq_iq_i^T$，其中（特征值）$\lambda_i\ge0,\ i=1,\cdots,n$。于是，我们有
\[
  \operatorname{tr}(YX)=\operatorname{tr}\left(Y\sum_{i=1}^n\lambda_iq_iq_i^T\right)
  =\sum_{i=1}^n\lambda_iq_i^TYq_i\ge0.
\]
以上表明 $Y\in(\mathbf{S}_+^n)^*$。

\noindent\textbf{例 2.25}\quad 范数锥的对偶。令 $\|\cdot\|$ 为定义在 $\mathbb{R}^n$ 上的范数。与之相关的锥 $K=\{(x,t)\in\mathbb{R}^{n+1}\mid\|x\|\le t\}$ 的对偶锥由其对偶范数定义，
\[
  K^*=\{(u,v)\in\mathbb{R}^{n+1}\mid \|u\|_*\le v\},
\]
这里的对偶范数由 $\|u\|_*=\sup\{u^Tx\mid\|x\|\le1\}$ 给出（参见 (A.1.6)）。

为证明这个结论，我们需要说明
\begin{equation}
  x^Tu+tv\ge0\quad \text{只要}\quad \|x\|\le t
  \Longleftrightarrow \|u\|_*\le v.
  \tag{2.20}
\end{equation}
首先，证明由右端关于 $(u,v)$ 的条件可以导出左端的条件。设 $\|u\|_*\le v$，并且对于一些 $t>0$ 有 $\|x\|\le t$。（如果 $t=0$，$x$ 必须是零，因此显然有 $u^Tx+vt\ge0$。）根据对偶锥的定义以及 $\|-x/t\|\le1$，我们有
\[
  u^T(-x/t)\le\|u\|_*\le v,
\]
因此，$u^Tx+vt\ge0$。

其次，我们证明 (2.20) 左端的条件可以导出 (2.20) 右端的条件。假设 $\|u\|_*>v$，即右端不成立。那么，根据对偶锥的定义，存在 $x$ 满足 $\|x\|\le1$ 及 $x^Tu>v$。取 $t=1$，我们有
\[
  u^T(-x)+v<0,
\]
这与 (2.20) 的左端相矛盾。

对偶锥满足一些性质，例如
\begin{itemize}
  \item $K^*$ 是闭凸锥。
  \item $K_1\subseteq K_2$ 可导出 $K_2^*\subseteq K_1^*$。
  \item 如果 $K$ 有非空内部，那么 $K^*$ 是尖的。
  \item 如果 $K$ 的闭包是尖的，那么 $K^*$ 有非空内部。
  \item $K^{**}$ 是 $K$ 的凸包的闭包。（因此，如果 $K$ 是凸和闭的，则 $K^{**}=K$。）（参见习题 2.31。）这些性质表明如果 $K$ 是一个正常锥，那么它的对偶 $K^*$ 也是；进一步地，有 $K^{**}=K$。
\end{itemize}

\subsection*{2.6.2\quad 广义不等式的对偶}

现在假设凸锥 $K$ 是正常锥，因此它可以导出一个广义不等式 $\preceq_K$。其对偶锥 $K^*$ 也是正常的，所以也能导出一个广义不等式。我们称广义不等式 $\preceq_{K^*}$ 为广义不等式 $\preceq_K$ 的对偶。

关于广义不等式及其对偶有一些重要的性质
\begin{itemize}
  \item $x\preceq_K y$ 当且仅当对于任意 $\lambda\succeq_{K^*}0$ 有 $\lambda^Tx\le\lambda^Ty$。
  \item $x\prec_K y$ 当且仅当对于任意 $\lambda\succeq_{K^*}0$ 和 $\lambda\ne0$，有 $\lambda^Tx<\lambda^Ty$。
\end{itemize}
因为 $K=K^{**}$，与 $\preceq_{K^*}$ 相关的对偶广义不等式为 $\preceq_K$，因此交换广义不等式及其对偶后，这些性质依然成立。作为一个具体的例子，我们可知 $\lambda\preceq_{K^*}\mu$ 的充要条件是对于所有 $x\succeq_K0$ 有 $\lambda^Tx\le\mu^Tx$。

\noindent\textbf{例 2.26}\quad 线性严格广义不等式的择一定理。设 $K\subseteq\mathbb{R}^m$ 为正常锥。考虑严格广义不等式
\begin{equation}
  Ax\prec_K b,
  \tag{2.21}
\end{equation}
其中 $x\in\mathbb{R}^n$。

我们将导出对于这个不等式的择一定理。假设它是不可行的，即仿射集合 $\{b-Ax\mid x\in\mathbb{R}^n\}$ 与开凸集 $\mathbf{int}\,K$ 不相交。那么存在一个分离超平面，即非零的 $\lambda\in\mathbb{R}^m$ 和 $\mu\in\mathbb{R}$ 使得对于任意 $x$ 有 $\lambda^T(b-Ax)\le\mu$，对于任意 $y\in\mathbf{int}\,K$ 有 $\lambda^Ty\ge\mu$。第一个条件表明 $A^T\lambda=0$ 及 $\lambda^Tb\le\mu$。第二个条件表明对于任意 $y\in K$ 有 $\lambda^Ty\ge\mu$，这种情况仅当 $\lambda\in K^*$ 和 $\mu\le0$ 时才可能发生。

综上，我们可知当 (2.21) 不可行时，存在 $\lambda$ 使得
\begin{equation}
  \lambda\ne0,\quad \lambda\succeq_{K^*}0,\quad A^T\lambda=0,\quad \lambda^Tb\le0.
  \tag{2.22}
\end{equation}
现在我们证明反方向，即如果 (2.22) 成立，那么，不等式组 (2.21) 不可能可行。假设不等式均成立，因为 $\lambda\ne0,\lambda\succeq_{K^*}0$ 及 $b-Ax\succ_K0$，我们有 $\lambda^T(b-Ax)>0$。但是根据 $A^T\lambda=0$，我们可以找到 $\lambda^T(b-Ax)=\lambda^Tb\le0$，而这是一个矛盾。

因此，不等式组 (2.21) 和 (2.22) 构成一对择一：对于任意 $A,b$，它们中仅有一个是可行的。（这是 (2.17),(2.18) 择一定理的推广；(2.17),(2.18) 是 $K=\mathbb{R}_+^m$ 时的特殊情况。）

\subsection*{2.6.3\quad 对偶不等式定义的最小元和极小元}

可以利用对偶广义不等式来刻画集合 $S\subseteq\mathbb{R}^m$（可能非凸）关于正常锥 $K$ 导出的广义不等式的最小元和极小元。

\noindent\textbf{最小元的对偶性质}

我们首先考虑最小元的性质。$x$ 是 $S$ 上关于广义不等式 $\preceq_K$ 的最小元的充要条件是，对于所有 $\lambda\succ_{K^*}0$，$x$ 是在 $z\in S$ 上极小化 $\lambda^Tz$ 的唯一最优解。几何上看，这意味着对于任意 $\lambda\succ_{K^*}0$，超平面
\[
  \{z\mid \lambda^T(z-x)=0\}
\]
是在 $x$ 处对 $S$ 的一个严格支撑超平面。（我们用严格支撑超平面表明这个超平面与 $S$ 只相交于 $x$。）注意此处并不要求 $S$ 是凸集。这可由图2.23来表示。

为说明这个结论，设 $x$ 是 $S$ 的最小元，即对于任意 $z\in S$ 有 $x\preceq_K z$，同时，令 $\lambda\succ_{K^*}0$，而 $z\in S,\ z\ne x$。因为 $x$ 是 $S$ 上的最小元，我们有 $z-x\succeq_K0$。根据 $\lambda\succ_{K^*}0$ 及 $z-x\succeq_K0,\ z-x\ne0$，可以得到 $\lambda^T(z-x)>0$。因为 $z$ 是 $S$ 上任意一个不等于 $x$ 的元素，所以 $x$ 是在 $z\in S$ 上极小化 $\lambda^Tz$ 的唯一解。反之，假设对于所有 $\lambda\succ_{K^*}0$，$x$ 是 $z\in S$ 上极小化 $\lambda^Tz$ 的唯一解，但 $x$ 不是 $S$ 的最小元。那么存在 $z\in S$ 满足 $z\not\succeq_K x$。因为 $z-x\not\succeq_K0$，存在 $\bar{\lambda}\succeq_{K^*}0$ 并且 $\bar{\lambda}^T(z-x)<0$。因此，对于 $\lambda\succ_{K^*}0$，在 $\bar{\lambda}$ 的邻域内有 $\lambda^T(z-x)<0$。这与 $x$ 是 $S$ 上极小化 $\lambda^Tz$ 的唯一解相矛盾。

\noindent\textbf{极小元的对偶性质}

现在我们转而讨论极小元的类似性质。此时，在必要和充分条件间存在一定的间隙。如果 $\lambda\succ_{K^*}0$ 并且 $x$ 在 $z\in S$ 上极小化 $\lambda^Tz$，那么 $x$ 是极小的，如图2.24所示。

为说明这点，假设 $\lambda\succ_{K^*}0$ 并且 $x$ 在 $S$ 上极小化 $\lambda^Tz$，但 $x$ 不是极小元，即存在 $z\in S$ 满足 $z\ne x,\ z\preceq_K x$。那么 $\lambda^T(x-z)>0$。这与我们的假设，即 $x$ 在 $S$ 上极小化了 $\lambda^Tz$，相矛盾。

其逆命题在一般情况下不成立：$S$ 上的极小元 $x$ 可以对于任何 $\lambda$ 都不是 $z\in S$ 上极小化 $\lambda^Tz$ 的解，如图2.25所示。此图表明了凸性在这个逆定理中的重要作用，当凸性成立时，逆定理是成立的。假设集合 $S$ 是凸集，可以说对于任意极小元 $x$，存在非零的 $\lambda\succeq_{K^*}0$ 使得 $x$ 在 $z\in S$ 上极小化 $\lambda^Tz$。

为证明这一点，假设 $x$ 是极小的，也就是说
\[
  ((x-K)\setminus\{x\})\cap S=\emptyset.
\]
对凸集 $(x-K)\setminus\{x\}$ 和 $S$ 应用超平面分离定理，我们可以得出：存在 $\lambda\ne0$ 和 $\mu$，使得对于所有 $y\in K$ 有 $\lambda^T(x-y)\le\mu$，对于所有 $z\in S$ 有 $\lambda^Tz\ge\mu$。根据第一个不等式，我们可知 $\lambda\succeq_{K^*}0$。由于 $x\in S$ 和 $x\in x-K$，我们有 $\lambda^Tx=\mu$，所以第二个不等式表明 $\mu$ 是 $S$ 上 $\lambda^Tz$ 的最小值。因此，$x$ 是 $S$ 上极小化 $\lambda^Tz$ 的一个解，这里 $\lambda\ne0,\lambda\succeq_{K^*}0$。

这个逆定理无法加强为 $\lambda\succ_{K^*}0$。反例表明，凸集 $S$ 上的极小元 $x$，可以对于任意 $\lambda\succ_{K^*}0$，都不是 $z\in S$ 中极小化 $\lambda^Tz$ 的解（参见图2.26，左图）。同时，并不是对于任意 $\lambda\succeq_{K^*}0$，在 $z\in S$ 上极小化 $\lambda^Tz$ 的解都一定是极小的（参见图2.26，右图）。

\noindent\textbf{例 2.27}\quad Pareto 最优制造前沿。我们考虑安排一个产品的生产，该产品需要 $n$ 种资源（例如劳动力、电、天然气、水）。这种产品可由多种方式进行制造。我们用资源向量 $x\in\mathbb{R}^n$ 表示各种制造方法，其中 $x_i$ 表示相应方法制造产品时消耗资源 $i$ 的数量。我们假设 $x_i\ge0$（即制造过程消耗资源）并且资源是有价值的（因此希望尽量少地消耗资源）。

\noindent\textbf{生产集合}\quad $P\subseteq\mathbb{R}^n$ 定义为表示制造方法的所有资源向量 $x$ 的集合。

$P$ 上的极小元（在分量不等式意义下）对应的制造方法称为 Pareto 最优，或有效的。$P$ 的极小元构成的集合称为有效制造前沿。

我们可以给出 Pareto 最优性的一个简单示例。我们称与资源向量 $x$ 相关的制造方法比与资源向量 $y$ 相关的制造方法更好，如果对于所有 $i$，$x_i\le y_i$ 并且存在某些 $i$，$x_i<y_i$。换言之，如果对于每一种资源，一个制造方法都不比另一个消耗得多，并且至少有一种资源，该方法确实消耗得更少，那么称这个方法比另一方法更优。也就是说 $x\preceq y,\ x\ne y$。那么，我们可以说一个制造方法是 Pareto 最优或有效的，如果不存在更好的制造方法。

我们可以通过在制造向量集合 $P$ 上对任意满足 $\lambda\succ0$ 的 $\lambda$ 极小化
\[
  \lambda^Tx=\lambda_1x_1+\cdots+\lambda_nx_n
\]
来得到 Pareto 最优制造方法（即极小的资源向量）。

这里的向量 $\lambda$ 有一个简单的解释：$\lambda_i$ 是资源 $i$ 的价格。通过在 $P$ 上极小化 $\lambda^Tx$，可以寻找最为便宜的制造方法（对于资源价格 $\lambda_i$）。只要价格是正的，其得到的制造方法可以保证是有效的。

图2.27展示了这些想法。

\section*{参考文献}

Minkowski 被认为第一个对凸集进行了系统的分析并且引入了很多基础性概念，例如支撑超平面，支撑超平面定理，Minkowski 距离函数（参见习题 3.34），凸集的极限点等。一些早期的综述可见 Bonnessen 和 Fenchel 的 [BF48]，Eggleston 的 [Egg58]，Klee 的 [Kle63]，以及 Valentinie 的 [Val64]。最近的一些书籍专注于凸集的几何特征，例如 Lay 的 [Lay82] 以及 Webster 的 [Web94]。Klee 的 [Kle71]，Fenchel 的 [Fen83]，Tikhomorov 的 [Tik90] 以及 Berger 的 [Ber90] 给出了非常有趣的综述，讲述了凸性的历史及其在数学方面的应用。

与线性规划问题相关联，对线性不等式及多项式集合已经有了充分的讨论。对于这点，我们在第 4 章末给出了相关的文献。关于线性不等式和线性规划的一些里程碑式的文献有：Motzkin 的 [Mot33]，von Neumann 和 Morgenstern 的 [vNM53]，Kantorovich 的 [Kan60]，Koopmans 的 [Koo51]，以及 Dantzig 的 [Dan63]。Dantzig 在 [Dan63，第 2 章] 中讨论了线性不等式，包括直至 1963 年前后的历史调研。

20世纪60年代中在非线性规划的研究中提出了广义不等式（参见 Luenberger 的 [Lue69，§8.2] 以及 Isii 的 [Isi64]），并且被扩展到锥优化问题中（参见第 4 章中的参考文献）。Bellman 和 Fan 的 [BF63] 是关于广义线性不等式集合（关于半正定锥）的一篇早期文献。

对于超平面分离定理证明的推广，我们推荐读者参看 Rockafellar 的 [Roc70, part III]，以及 Hiriart-Urruty 和 Lemaréchal 的 [HUL93, volume 1，§III4]。Dantzig 的 [Dan63，第 21 页] 包含了 von Neumann 和 Morgenstern 在 [vNM53，第 138 页] 给出的择一定理。关于择一定理的更多文献，参见第 5 章。

例 2.27 中的术语（包括 Pareto 最优性，有效制造，价格 $\lambda$ 的解释）在 Luenberger 的 [Lue95] 中进行了详尽的讨论。

凸的几何性质在经典的力矩理论（Krein 和 Nudelman 的 [KN77]，Karlin 和 Studden 的 [KS66]）中有显著的作用。一个著名的例子是非负多项式与力矩锥的对偶性，参见习题 2.37。

\section*{习题}

\noindent\textbf{凸性定义}

\noindent\textbf{2.1}\quad 设 $C\subseteq\mathbb{R}^n$ 为一个凸集且 $x_1,\cdots,x_k\in C$。令 $\theta_1,\cdots,\theta_k\in\mathbb{R}$ 满足 $\theta_i\ge0,\ \theta_1+\cdots+\theta_k=1$。证明 $\theta_1x_1+\cdots+\theta_kx_k\in C$。（凸性的定义是指此式在 $k=2$ 时成立；你需要证明对任意 $k$ 的情况。）提示：对 $k$ 进行归纳。

\noindent\textbf{2.2}\quad 证明一个集合是凸集当且仅当它与任意直线的交是凸的。证明一个集合是仿射的，当且仅当它与任意直线的交是仿射的。

\noindent\textbf{2.3}\quad 中点凸性。集合 $C$ 是中点凸的，当 $C$ 中任意两点 $a,b$ 的平均或中点 $(a+b)/2$ 也属于 $C$。显然凸集是中点凸的。可以证明在一些很微弱的条件下，中点凸可以导出凸性。作为一个简单的例子，证明如果 $C$ 是闭和中点凸的，那么 $C$ 是凸集。

\noindent\textbf{2.4}\quad 证明集合 $S$ 的凸包是所有包含 $S$ 的凸集的交。（同样的方法可以用来证明集合 $S$ 的锥包、仿射包或线性包分别是所有包含 $S$ 的锥集合、仿射集合或半空间的交集。）

\noindent\textbf{例子}

\noindent\textbf{2.5}\quad 两个平行的超平面 $\{x\in\mathbb{R}^n\mid a^Tx=b_1\}$ 和 $\{x\in\mathbb{R}^n\mid a^Tx=b_2\}$ 之间的距离是多少？

\noindent\textbf{2.6}\quad 什么时候超平面包含另一个？给出使下式成立的条件
\[
  \{x\mid a^Tx\le b\}\subseteq\{x\mid \tilde{a}^Tx\le \tilde{b}\}
\]
（其中 $a\ne0,\tilde{a}\ne0$）。另外，找出使得两个半空间相等的条件。

\noindent\textbf{2.7}\quad 半空间的 Voronoi 描述。令 $a$ 和 $b$ 为 $\mathbb{R}^n$ 上互异的两点。证明所有距离 $a$ 比距离 $b$ 近（Euclid 范数下）的点的集合，即 $\{x\mid \|x-a\|_2\le\|x-b\|_2\}$，是一个超平面。用形如 $c^Tx\le d$ 的不等式进行显式表示并给出图像。

\noindent\textbf{2.8}\quad 下面的集合，$S$ 中哪些是多面体？如果可能，将 $S$ 表示为 $S=\{x\mid Ax\preceq b,\ Fx=g\}$ 的形式。
\[
  \begin{array}{ll}
  \text{(a)} & S=\{y_1a_1+y_2a_2\mid -1\le y_1\le1,\ -1\le y_2\le1\},\quad \text{其中 }a_1,a_2\in\mathbb{R}^n.\\[0.4em]
  \text{(b)} & S=\{x\in\mathbb{R}^n\mid x\succeq0,\ \mathbf{1}^Tx=1,\ \sum_{i=1}^n x_ia_i=b_1,\ \sum_{i=1}^n x_ia_i^2=b_2\},\\
  & \text{其中 }a_1,\cdots,a_n\in\mathbb{R}\text{ 并且 }b_1,b_2\in\mathbb{R}.\\[0.4em]
  \text{(c)} & S=\{x\in\mathbb{R}^n\mid x\succeq0,\ x^Ty\le1\text{ 对于所有满足 }\|y\|_2=1\text{ 的 }y\}.\\[0.4em]
  \text{(d)} & S=\{x\in\mathbb{R}^n\mid x\succeq0,\ x^Ty\le1\text{ 对于所有满足 }\sum_{i=1}^n |y_i|=1\text{ 的 }y\}.
  \end{array}
\]

\noindent\textbf{2.9}\quad Voronoi 集合与多面体分解。令 $x_0,\cdots,x_K\in\mathbb{R}^n$。考虑所有距离 $x_0$ 比距离其他 $x_i$ 点更近（Euclid 范数下）的点组成的集合，即
\[
  V=\{x\in\mathbb{R}^n\mid \|x-x_0\|_2\le\|x-x_i\|_2,\ i=1,\cdots,K\}.
\]
$V$ 称为围绕 $x_0$ 的关于 $x_1,\cdots,x_K$ 的 Voronoi 区域。

\noindent(a) 证明 $V$ 是一个多面体，并将 $V$ 表示为 $V=\{x\mid Ax\preceq b\}$ 的形式。

\noindent(b) 反之，给定一个内部非空的多面体 $P$，说明如何寻找 $x_0,\cdots,x_K$ 使得这个多面体是 $x_0$ 关于 $x_1,\cdots,x_K$ 的 Voronoi 区域。

\noindent(c) 我们也可考虑集合
\[
  V_k=\{x\in\mathbb{R}^n\mid \|x-x_k\|_2\le\|x-x_i\|_2,\ i\ne k\}.
\]
集合 $V_k$ 由 $\mathbb{R}^n$ 上的点构成，并且距离集合 $\{x_0,\cdots,x_K\}$ 最近的点是 $x_k$。集合 $V_0,\cdots,V_K$ 给出了 $\mathbb{R}^n$ 的多面体分解。更准确地，集合 $V_k$ 是多面体，$\bigcup_{k=0}^K V_k=\mathbb{R}^n$，并且对于任意 $i\ne j$ 有 $\mathbf{int}\,V_i\cap\mathbf{int}\,V_j=\emptyset$，即 $V_i$ 与 $V_j$ 最多在边界相交。假设 $P_1,\cdots,P_m$ 为多面体，并且使得 $\bigcup_{i=1}^m P_i=\mathbb{R}^n$，对于 $i\ne j$ 有 $\mathbf{int}\,P_i\cap\mathbf{int}\,P_j=\emptyset$。这种对于 $\mathbb{R}^n$ 的多面体分解是否可以用 Voronoi 区域来表示，这些区域由适当的点集产生。

\noindent\textbf{2.10}\quad 二次不等式的解集。令 $C\subseteq\mathbb{R}^n$ 为下列二次不等式的解集，
\[
  C=\{x\in\mathbb{R}^n\mid x^TAx+b^Tx+c\le0\},
\]
其中 $A\in\mathbf{S}^n,\ b\in\mathbb{R}^n,\ c\in\mathbb{R}$。

\noindent(a) 证明：如果 $A\succeq0$，那么 $C$ 是凸集。

\noindent(b) 证明：如果对某些 $\lambda\in\mathbb{R}$ 有 $A+\lambda gg^T\succeq0$，那么 $C$ 和由 $g^Tx+h=0$（这里 $g\ne0$）定义的超平面的交集是凸集。

以上命题的逆命题是否成立？

\noindent\textbf{2.11}\quad 双曲集合。证明双曲集合 $\{x\in\mathbb{R}_+^2\mid x_1x_2\ge1\}$ 是凸集。更一般地，证明 $\{x\in\mathbb{R}_+^n\mid \prod_{i=1}^n x_i\ge1\}$ 是凸的。提示：如果 $a,b\ge0$ 并且 $0\le\theta\le1$，那么 $a^\theta b^{1-\theta}\le\theta a+(1-\theta)b$；参见 \S3.1.9。

\noindent\textbf{2.12}\quad 下面的集合哪些是凸集？
\[
  \begin{array}{ll}
  \text{(a)} & \text{平板，即形如 }\{x\in\mathbb{R}^n\mid \alpha\le a^Tx\le\beta\}\text{ 的集合。}\\[0.3em]
  \text{(b)} & \text{矩形，即形如 }\{x\in\mathbb{R}^n\mid \alpha_i\le x_i\le\beta_i,\ i=1,\cdots,n\}\text{ 的集合。}\\
  & \text{当 }n>2\text{ 时，矩形有时也称为超矩形。}\\[0.3em]
  \text{(c)} & \text{楔形，即 }\{x\in\mathbb{R}^n\mid a_1^Tx\le b_1,\ a_2^Tx\le b_2\}.\\[0.3em]
  \text{(d)} & \text{距离给定点比距离给定集合近的点构成的集合，即}\\
  & \{x\mid \|x-x_0\|_2\le\|x-y\|_2,\ \forall y\in S\},\quad \text{其中 }S\subseteq\mathbb{R}^n.\\[0.3em]
  \text{(e)} & \text{距离一个集合比另一个集合更近的点的集合，即}\\
  & \{x\mid \mathbf{dist}(x,S)\le\mathbf{dist}(x,T)\},\\
  & \text{其中 }S,T\subseteq\mathbb{R}^n,\quad
    \mathbf{dist}(x,S)=\inf\{\|x-z\|_2\mid z\in S\}.\\[0.3em]
  \text{(f)} & \text{[HUL93，第 1 卷，第 93 页] 集合 }\{x\mid x+S_2\subseteq S_1\},\\
  & \text{其中 }S_1,S_2\subseteq\mathbb{R}^n\text{ 并且 }S_1\text{ 是凸集。}\\[0.3em]
  \text{(g)} & \text{到 }a\text{ 的距离与到 }b\text{ 的距离之比不超过到某一固定分数 }\theta\text{ 的点的集合，}\\
  & \text{即集合 }\{x\mid \|x-a\|_2\le\theta\|x-b\|_2\}\text{。可以假设 }a\ne b\text{ 及 }0\le\theta\le1.
  \end{array}
\]

\noindent\textbf{2.13}\quad 外积的闭包。考虑秩为 $k$ 的外积，即 $\{XX^T\mid X\in\mathbb{R}^{n\times k},\ \operatorname{rank}X=k\}$。试用简单的形式描述其闭包。

\noindent\textbf{2.14}\quad 扩展和限制集合。令 $S\subseteq\mathbb{R}^n$，用 $\|\cdot\|$ 表示 $\mathbb{R}^n$ 上的范数。

\noindent(a) 对于 $a\ge0$ 我们定义 $S_a$ 为 $\{x\mid \mathbf{dist}(x,S)\le a\}$，其中 $\mathbf{dist}(x,S)=\inf_{y\in S}\|x-y\|$。我们称 $S_a$ 为 $S$ 的扩展或延伸，其幅度为 $a$。证明如果 $S$ 是凸集，那么 $S_a$ 是凸的。

\noindent(b) 对于 $a\ge0$，我们定义 $S_{-a}=\{x\mid B(x,a)\subseteq S\}$，其中 $B(x,a)$ 是以 $x$ 为中心、$a$ 为半径的球（在范数 $\|\cdot\|$ 意义下）。我们称 $S_{-a}$ 为 $S$ 的收缩或限制，其幅度为 $a$，因为 $S_{-a}$ 是由所有离 $\mathbb{R}^n\setminus S$ 的距离至少为 $a$ 的点的集合。证明如果 $S$ 是凸集，那么 $S_{-a}$ 也是凸集。

\noindent\textbf{2.15}\quad 一些概率分布集合。令 $x$ 为服从分布 $\operatorname{prob}(x=a_i)=p_i,\ i=1,\cdots,n$ 的实数随机变量，其中 $a_1<a_2<\cdots<a_n$。当然 $p\in\mathbb{R}^n$ 在一个标准概率单纯形 $P=\{p\mid \mathbf{1}^Tp=1,\ p\succeq0\}$ 中。下面哪些条件在 $p$ 中是凸的？（即满足下面哪些条件的 $p\in P$ 的集合是凸集？）
\[
  \begin{array}{ll}
  \text{(a)} & \alpha\le \mathbf{E}f(x)\le\beta,\quad
    \text{其中 }\mathbf{E}f(x)\text{ 为 }f(x)\text{ 的期望，即 }\mathbf{E}f(x)=\sum_{i=1}^n p_if(a_i)\\
  & \text{（给定函数 }f:\mathbb{R}\to\mathbb{R}\text{）。}\\[0.3em]
  \text{(b)} & \operatorname{prob}(x>\alpha)\le\beta.\\
  \text{(c)} & \mathbf{E}|x^3|\le\alpha\mathbf{E}|x|.\\
  \text{(d)} & \mathbf{E}x^2\le\alpha.\\
  \text{(e)} & \mathbf{E}x^2\ge\alpha.\\
  \text{(f)} & \operatorname{var}(x)\le\alpha,\quad \text{其中 }\operatorname{var}(x)=\mathbf{E}(x-\mathbf{E}x)^2\text{ 为 }x\text{ 的方差。}\\
  \text{(g)} & \operatorname{var}(x)\ge\alpha.\\
  \text{(h)} & \operatorname{quartile}(x)\ge\alpha,\quad \text{其中 }\operatorname{quartile}(x)=\inf\{\beta\mid \operatorname{prob}(x\le\beta)\ge0.25\}.\\
  \text{(i)} & \operatorname{quartile}(x)\le\alpha.
  \end{array}
\]

\noindent\textbf{保凸运算}

\noindent\textbf{2.16}\quad 证明如果 $S_1$ 和 $S_2$ 是 $\mathbb{R}^{m+n}$ 中的凸集，那么它们的部分和
\[
  S=\{(x,y_1+y_2)\mid x\in\mathbb{R}^m,\ y_1,y_2\in\mathbb{R}^n,\ (x,y_1)\in S_1,\ (x,y_2)\in S_2\}
\]
也是凸的。

\noindent\textbf{2.17}\quad 透视函数下的多面体集合。在这个问题中，我们研究超平面、半空间及多面体在透视函数 $P(x,t)=x/t$ 下的像，其中 $\operatorname{dom}P=\mathbb{R}^n\times\mathbb{R}_{++}$。对于下面每个集合 $C$，给出形如
\[
  P(C)=\{v/t\mid (v,t)\in C,\ t>0\}
\]
的简单表示。
\[
  \begin{array}{ll}
  \text{(a)} & \text{多面体 }C=\mathbf{conv}\{(v_1,t_1),\cdots,(v_K,t_K)\},\quad \text{其中 }v_i\in\mathbb{R}^n,\ t_i>0.\\[0.3em]
  \text{(b)} & \text{超平面 }C=\{(v,t)\mid f^Tv+gt=h\}\quad \text{（其中 }f\text{ 和 }g\text{ 不均为零）。}\\[0.3em]
  \text{(c)} & \text{半空间 }C=\{(v,t)\mid f^Tv+gt\le h\}\quad \text{（其中 }f\text{ 和 }g\text{ 不均为零）。}\\[0.3em]
  \text{(d)} & \text{多面体 }C=\{(v,t)\mid Fv+gt\preceq h\}.
  \end{array}
\]

\noindent\textbf{2.18}\quad 可逆的线性分式函数。令 $f:\mathbb{R}^n\to\mathbb{R}^n$ 为线性分式函数
\[
  f(x)=(Ax+b)/(c^Tx+d),\qquad \operatorname{dom}f=\{x\mid c^Tx+d>0\}.
\]
设矩阵
\[
  Q=\begin{bmatrix} A & b\\ c^T & d \end{bmatrix}
\]
非奇异。证明 $f$ 可逆并且 $f^{-1}$ 也是一个线性分式映射。利用 $A,b,c$ 和 $d$ 显式地给出 $f^{-1}$ 及其定义域的表达式。提示：用 $Q$ 来表示 $f^{-1}$ 可能会更容易些。

\noindent\textbf{2.19}\quad 线性分式函数和凸集。令 $f:\mathbb{R}^m\to\mathbb{R}^n$ 为线性分式函数
\[
  f(x)=(Ax+b)/(c^Tx+d),\qquad \operatorname{dom}f=\{x\mid c^Tx+d>0\}.
\]
在这个问题中，我们研究凸集 $C$ 在 $f$ 下的原像，即
\[
  f^{-1}(C)=\{x\in\operatorname{dom}f\mid f(x)\in C\}.
\]
对下面的每个集合 $C\subseteq\mathbb{R}^n$，给出 $f^{-1}(C)$ 的简单描述。
\[
  \begin{array}{ll}
  \text{(a)} & \text{半空间 }C=\{y\mid g^Ty\le h\}\quad \text{（其中 }g\ne0\text{）。}\\
  \text{(b)} & \text{多面体 }C=\{y\mid Gy\preceq h\}.\\
  \text{(c)} & \text{椭球 }\{y\mid y^TP^{-1}y\le1\}\quad \text{（其中 }P\in\mathbf{S}_{++}^n\text{）。}\\
  \text{(d)} & \text{线性矩阵不等式的解集 }C=\{y\mid y_1A_1+\cdots+y_nA_n\preceq B\},\\
  & \text{其中 }A_1,\cdots,A_n,B\in\mathbf{S}^p.
  \end{array}
\]

\noindent\textbf{分离定理与支撑超平面}

\noindent\textbf{2.20}\quad 线性方程组的严格正解。设 $A\in\mathbb{R}^{m\times n},\ b\in\mathbb{R}^m$，其中 $b\in\mathcal{R}(A)$。证明存在 $x$ 满足
\[
  x\succ0,\qquad Ax=b
\]
的充要条件是不存在 $\lambda$ 满足
\[
  A^T\lambda\succeq0,\qquad A^T\lambda\ne0,\qquad b^T\lambda\le0.
\]
提示：首先由线性代数证明：对于所有满足 $Ax=b$ 的 $x$，$c^Tx=d$ 的充要条件是存在向量 $\lambda$ 满足 $c=A^T\lambda,\ d=b^T\lambda$。

\noindent\textbf{2.21}\quad 分离超平面的集合。设 $C$ 和 $D$ 为 $\mathbb{R}^n$ 的不相交的子集。考虑集合 $(a,b)\in\mathbb{R}^{n+1}$，它满足对任意 $x\in C$ 有 $a^Tx\le b$；对任意 $x\in D$ 有 $a^Tx\ge b$。证明这个集合是一个凸锥（并且如果没有分离 $C$ 和 $D$ 的超平面，那么它是单点集 $\{0\}$）。

\noindent\textbf{2.22}\quad 完成 \S2.5.1 中超平面分离定理的证明：证明对于两个不相交的凸集 $C$ 和 $D$ 存在分离超平面。可以利用 \S2.5.1 中已证的结论，即当在两个集合中存在点，其距离与两个集合间的距离相等时，存在分离超平面。

提示：如果 $C$ 和 $D$ 是不相交的凸集，那么集合 $\{x-y\mid x\in C,\ y\in D\}$ 是凸集并且不包含原点。

\noindent\textbf{2.23}\quad 给出两个不相交的闭凸集不能被严格分离的例子。

\noindent\textbf{2.24}\quad 支撑超平面。

\noindent(a) 将闭凸集 $\{x\in\mathbb{R}_+^2\mid x_1x_2\ge1\}$ 表示为半空间的交集。

\noindent(b) 令 $C=\{x\in\mathbb{R}^n\mid \|x\|_\infty\le1\}$ 表示 $\mathbb{R}^n$ 空间中的单位 $\ell_\infty$-范数球，并令 $\hat{x}$ 为 $C$ 的边界上的点。显式地写出集合 $C$ 在 $\hat{x}$ 处的支撑超平面。

\noindent\textbf{2.25}\quad 内部和外部多面体逼近。令 $C\subseteq\mathbb{R}^n$ 为闭凸集，并设 $x_1,\cdots,x_K$ 在 $C$ 的边界上。设对于每个 $i$，$a_i^T(x-x_i)=0$ 定义了 $C$ 在 $x_i$ 处的一个支撑超平面，即 $C\subseteq\{x\mid a_i^T(x-x_i)\le0\}$。考虑两个多面体
\[
  P_{\mathrm{inner}}=\mathbf{conv}\{x_1,\cdots,x_K\},\qquad
  P_{\mathrm{outer}}=\{x\mid a_i^T(x-x_i)\le0,\ i=1,\cdots,K\}.
\]
证明 $P_{\mathrm{inner}}\subseteq C\subseteq P_{\mathrm{outer}}$ 并画出图像进行说明。

\noindent\textbf{2.26}\quad 支撑函数。集合 $C\subseteq\mathbb{R}^n$ 的支撑函数定义为
\[
  S_C(y)=\sup\{y^Tx\mid x\in C\}.
\]
（我们允许 $S_C(y)$ 取值为 $+\infty$。）设 $C$ 和 $D$ 是 $\mathbb{R}^n$ 上的闭凸集。证明 $C=D$ 当且仅当它们的支撑函数相等。

\noindent\textbf{2.27}\quad 支撑超平面定理的逆定理。设集合 $C$ 是闭的，含有非空内部并且在其边界上的每一点都有支撑超平面。证明 $C$ 是凸集。

\noindent\textbf{凸锥及广义不等式}

\noindent\textbf{2.28}\quad $n=1,2,3$ 时的半正定锥。对于 $n=1,2,3$，用矩阵系数和普通不等式给出半正定锥 $\mathbf{S}_+^n$ 的显式表示。为表示 $n=1,2,3$ 时 $\mathbf{S}^n$ 的一般元素，请用下面的符号
\[
  x_1,\qquad
  \begin{bmatrix} x_1 & x_2\\ x_2 & x_3 \end{bmatrix},\qquad
  \begin{bmatrix} x_1 & x_2 & x_3\\ x_2 & x_4 & x_5\\ x_3 & x_5 & x_6 \end{bmatrix}.
\]

\noindent\textbf{2.29}\quad $\mathbb{R}^2$ 中的锥。设 $K\subseteq\mathbb{R}^2$ 为一个闭凸锥。

\noindent(a) 给出 $K$ 在其元素的极坐标形式 $(x=r(\cos\phi,\sin\phi),\ r\ge0)$ 下的简单描述。

\noindent(b) 给出 $K^*$ 的简单描述，并绘图说明 $K$ 和 $K^*$ 之间的关系。

\noindent(c) 什么时候 $K$ 是尖的。

\noindent(d) 什么时候 $K$ 是正常的（因此，定义了广义不等式）？绘图说明当 $K$ 正则时，$x\preceq_K y$ 的含义。

\noindent\textbf{2.30}\quad 广义不等式的性质。证明 \S2.4.1 中所列的（不严格和严格的）广义不等式的性质。

\noindent\textbf{2.31}\quad 对偶锥的性质。令 $K^*$ 为凸锥 $K$ 的对偶锥，如 (2.19) 的定义。证明下面的性质
\[
  \begin{array}{ll}
  \text{(a)} & K^*\text{ 确实是凸锥。}\\
  \text{(b)} & K_1\subseteq K_2\text{ 表明 }K_2^*\subseteq K_1^*.\\
  \text{(c)} & K^*\text{ 是闭集。}\\
  \text{(d)} & K^*\text{ 的内部由 }\mathbf{int}\,K^*=\{y\mid y^Tx>0,\ \forall x\in\mathbf{cl}\,K\}\text{ 给出。}\\
  \text{(e)} & \text{如果 }K\text{ 具有非空内部，那么 }K^*\text{ 是尖的。}\\
  \text{(f)} & K^{**}\text{ 是 }K\text{ 的闭包。（因此，如果 }K\text{ 是闭的，那么 }K^{**}=K\text{。）}\\
  \text{(g)} & \text{如果 }K\text{ 的闭包是尖的，那么 }K^*\text{ 有非空内部。}
  \end{array}
\]

\noindent\textbf{2.32}\quad 寻找 $\{Ax\mid x\succeq0\}$ 的对偶锥，其中 $A\in\mathbb{R}^{m\times n}$。

\noindent\textbf{2.33}\quad 单调非负锥。我们定义单调非负锥为
\[
  K_{\mathrm{m}+}=\{x\in\mathbb{R}^n\mid x_1\ge x_2\ge\cdots\ge x_n\ge0\},
\]
即所有分量按非增排序的非负向量。

\noindent(a) 说明 $K_{\mathrm{m}+}$ 是正常锥。

\noindent(b) 找到对偶锥 $K_{\mathrm{m}+}^*$。提示：利用恒等式
\begin{align*}
  \sum_{i=1}^n x_iy_i
  &= (x_1-x_2)y_1+(x_2-x_3)(y_1+y_2)\\
  &\quad +(x_3-x_4)(y_1+y_2+y_3)+\cdots\\
  &\quad +(x_{n-1}-x_n)(y_1+\cdots+y_{n-1})+x_n(y_1+\cdots+y_n).
\end{align*}

\noindent\textbf{2.34}\quad 字典锥及排序。字典锥定义为
\[
  K_{\mathrm{lex}}=\{0\}\cup\{x\in\mathbb{R}^n\mid x_1=\cdots=x_k=0,\ x_{k+1}>0,\ \text{对某个 }k,\ 0\le k<n\},
\]
即所有第一个非零分量（如果存在）为正的向量。

\noindent(a) 验证 $K_{\mathrm{lex}}$ 是锥，但不是正常锥。

\noindent(b) 我们定义 $\mathbb{R}^n$ 中的字典排序如下：$x\preceq_{\mathrm{lex}}y$ 当且仅当 $y-x\in K_{\mathrm{lex}}$。（因为 $K_{\mathrm{lex}}$ 不是正常锥，字典排序不是广义不等式。）说明字典排序是一个线性序：对于任意 $x,y\in\mathbb{R}^n$，或者 $x\preceq_{\mathrm{lex}}y$ 或者 $y\preceq_{\mathrm{lex}}x$。所以，任意向量集合都可以关于字典锥进行排序，得到类似于应用于字典的排序。

\noindent(c) 寻找 $K_{\mathrm{lex}}^*$。

\noindent\textbf{2.35}\quad 谱正矩阵。矩阵 $X\in\mathbf{S}^n$ 称为谱正，如果对于所有 $z\succeq0$ 有 $z^TXz\ge0$。验证谱正矩阵的集合是一个正常锥，并寻找到其对偶锥。

\noindent\textbf{2.36}\quad Euclid 距离矩阵。令 $x_1,\cdots,x_n\in\mathbb{R}^k$。由 $D_{ij}=\|x_i-x_j\|_2^2$ 定义的矩阵 $D\in\mathbf{S}^n$ 称为 Euclid 距离矩阵。它满足一些显然的性质，例如 $D_{ij}=D_{ji}$，$D_{ii}=0$，$D_{ij}\ge0$，及（由三角不等式得到的）$D_{ik}^{1/2}\le D_{ij}^{1/2}+D_{jk}^{1/2}$。我们现在提出问题：什么时候一个矩阵 $D\in\mathbf{S}^n$ 是（关于某些 $k$ 的 $\mathbb{R}^k$ 空间中某些点的）Euclid 距离矩阵？一个著名的结果可以回答这个问题：$D\in\mathbf{S}^n$ 为 Euclid 矩阵的充要条件是，对于所有满足 $\mathbf{1}^Tx=0$ 的 $x$ 都有 $D_{ii}=0$ 及 $x^TDx\le0$（参见 \S8.3.3）。

证明 Euclid 距离矩阵集合是一个凸锥。

\noindent\textbf{2.37}\quad 非负多项式及 Hankel 线性矩阵不等式。令 $K_{\mathrm{pol}}$ 为 $\mathbb{R}$ 中 $2k$ 阶非负多项式的（系数）集合：
\[
  K_{\mathrm{pol}}=\{x\in\mathbb{R}^{2k+1}\mid x_1+x_2t+x_3t^2+\cdots+x_{2k+1}t^{2k}\ge0,\ \forall t\in\mathbb{R}\}.
\]

\noindent(a) 证明 $K_{\mathrm{pol}}$ 是一个正常锥。

\noindent(b) 一个基本的结果表明 $2k$ 阶多项式在 $\mathbb{R}$ 上非负的充要条件是，它可以被表示为两个 $k$ 或更低阶次的多项式的平方和。换言之，$x\in K_{\mathrm{pol}}$，当且仅当多项式
\[
  p(t)=x_1+x_2t+x_3t^2+\cdots+x_{2k+1}t^{2k}
\]
可以被表示为
\[
  p(t)=r(t)^2+s(t)^2,
\]
其中 $r$ 和 $s$ 为阶次为 $k$ 的多项式。

利用这个结果来说明
\[
  K_{\mathrm{pol}}=
  \left\{x\in\mathbb{R}^{2k+1}\ \middle|\ x_i=\sum_{m+n=i+1}Y_{mn}\text{ 对某些 }Y\in\mathbf{S}_+^{k+1}\right\}.
\]
也就是说，$p(t)=x_1+x_2t+x_3t^2+\cdots+x_{2k+1}t^{2k}$ 非负的充要条件是，存在矩阵 $Y\in\mathbf{S}_+^{k+1}$ 使得
\[
  x_1=Y_{11},\qquad
  x_2=Y_{12}+Y_{21},\qquad
  x_3=Y_{13}+Y_{22}+Y_{31},\qquad \cdots,\qquad
  x_{2k+1}=Y_{k+1,k+1}.
\]

\noindent(c) 证明 $K_{\mathrm{pol}}^*=K_{\mathrm{han}}$，其中
\[
  K_{\mathrm{han}}=\{z\in\mathbb{R}^{2k+1}\mid H(z)\succeq0\},
\]
而
\[
  H(z)=
  \begin{bmatrix}
  z_1 & z_2 & z_3 & \cdots & z_k & z_{k+1}\\
  z_2 & z_3 & z_4 & \cdots & z_{k+1} & z_{k+2}\\
  z_3 & z_4 & z_5 & \cdots & z_{k+2} & z_{k+3}\\
  \vdots & \vdots & \vdots & \ddots & \vdots & \vdots\\
  z_k & z_{k+1} & z_{k+2} & \cdots & z_{2k-1} & z_{2k}\\
  z_{k+1} & z_{k+2} & z_{k+3} & \cdots & z_{2k} & z_{2k+1}
  \end{bmatrix}.
\]
（这是关于系数 $z_1,\cdots,z_{2k+1}$ 的 Hankel 矩阵。）

\noindent(d) 令 $K_{\mathrm{mom}}$ 为所有具有 $(1,t,t^2,\cdots,t^{2k})$ 形式的向量所组成集合的锥包，其中 $t\in\mathbb{R}$。证明 $y\in K_{\mathrm{mom}}$，当且仅当 $y_1\ge0$ 并且对于某些随机变量 $u$，有
\[
  y=y_1(1,\mathbf{E}u,\mathbf{E}u^2,\cdots,\mathbf{E}u^{2k}).
\]
换言之，$K_{\mathrm{mom}}$ 的元素是 $\mathbb{R}$ 上所有可能扰动的矩向量的非负倍。证明 $K_{\mathrm{pol}}=K_{\mathrm{mom}}^*$。

\noindent(e) 结合 (c) 和 (d) 的结果，得出 $K_{\mathrm{pol}}=K_{\mathrm{mom}}^*$ 的结论。

作为一个说明 $K_{\mathrm{mom}}$ 和 $K_{\mathrm{han}}$ 之间关系的例子，取 $k=2,\ z=(1,0,0,0,1)$。证明 $z\in K_{\mathrm{han}},\ z\notin K_{\mathrm{mom}}$。找出 $K_{\mathrm{mom}}$ 中一个趋向于 $z$ 的显式点列。

\noindent\textbf{2.38}\quad [Roc70，第 15, 61 页] 从集合构造凸锥。

\noindent(a) 集合 $C$ 的障碍锥定义为使得 $y^Tx$ 在 $x\in C$ 中上有界的所有向量 $y$ 的集合。换言之，一个非零向量属于障碍锥，当且仅当它是某个包含 $C$ 的半空间 $\{x\mid y^Tx\le\alpha\}$ 的法向量。验证障碍锥是一个凸锥（不需要对 $C$ 的任何假设）。

\noindent(b) 集合 $C$ 的回收锥（也称为渐进锥）定义为对于每个 $x\in C$ 和所有 $t\ge0$ 都有 $x-ty\in C$ 的所有向量 $y$ 的集合。说明凸集的回收锥是凸锥。说明，如果 $C$ 是非空、闭和凸的，那么 $C$ 的回收锥是其障碍锥的对偶。

\noindent(c) 集合 $C$ 在其边界点 $x_0$ 处的法向锥定义为使得对于所有 $x\in C$ 都有 $y^T(x-x_0)\le0$ 的所有向量 $y$ 的集合（即在 $x_0$ 处定义了 $C$ 的支撑超平面的向量的集合）。说明法向锥是凸锥（不需要对 $C$ 的任何假设）。给出多面体 $\{x\mid Ax\preceq b\}$ 在其边界上一点的法向锥的简单描述。

\noindent\textbf{2.39}\quad 锥的分离。令 $K$ 和 $\tilde{K}$ 是两个内部非空且不相交的凸锥，即 $\mathbf{int}\,K\cap\mathbf{int}\,\tilde{K}=\emptyset$。证明存在非零 $y$ 使得 $y\in K^*,\ -y\in\tilde{K}^*$。

对于 $K=\tilde{K}$，这表明如果锥 $K$ 具有非空内部，那么 $K^*$ 是尖的。

\end{document}
